\documentclass[11pt]{article}

\usepackage[letterpaper,margin=0.82in,headheight=15pt]{geometry}
\usepackage[T1]{fontenc}
\usepackage{lmodern}
\usepackage{microtype}
\usepackage{amsmath,amssymb,amsthm,bm,mathtools}
\usepackage{booktabs,longtable,tabularx,array,multirow}
\usepackage{enumitem}
\usepackage{xcolor}
\usepackage{fancyhdr}
\usepackage{graphicx}
\usepackage{tikz-cd}
\usepackage{hyperref}
\usepackage[nameinlink,noabbrev]{cleveref}

\definecolor{navy}{HTML}{17324D}
\definecolor{teal}{HTML}{147D92}
\definecolor{orange}{HTML}{D47A27}
\definecolor{softblue}{HTML}{EEF4F7}
\definecolor{softgray}{HTML}{F3F5F6}
\definecolor{darkgray}{HTML}{46545E}
\definecolor{softorange}{HTML}{FFF4E8}

\hypersetup{
  colorlinks=true,
  linkcolor=navy,
  citecolor=teal,
  urlcolor=teal,
  pdftitle={Volume II: Common Nucleon GTMD Operators and Fock-Space Overlaps},
  pdfauthor={DeuteronWigner project}
}

\pagestyle{fancy}
\fancyhf{}
\fancyhead[L]{\small\color{darkgray}DeuteronWigner formalism - Volume II}
\fancyhead[R]{\small\color{darkgray}Common nucleon GTMD overlaps}
\fancyfoot[L]{\small\color{darkgray}28 July 2026}
\fancyfoot[R]{\small\color{darkgray}\thepage}
\renewcommand{\headrulewidth}{0.4pt}

\setlist[itemize]{leftmargin=1.5em,itemsep=2pt,topsep=3pt}
\setlist[enumerate]{leftmargin=1.7em,itemsep=2pt,topsep=3pt}
\setlength{\parindent}{0pt}
\setlength{\parskip}{5pt}
\setlength{\emergencystretch}{2em}
\renewcommand{\arraystretch}{1.16}
\allowdisplaybreaks

\newtheorem{definition}{Definition}[section]
\newtheorem{axiom}[definition]{Axiom}
\newtheorem{principle}[definition]{Design principle}
\newtheorem{requirement}[definition]{Requirement}
\newtheorem{proposition}[definition]{Proposition}
\newtheorem{theorem}[definition]{Theorem}
\newtheorem{lemma}[definition]{Lemma}
\newtheorem{corollary}[definition]{Corollary}
\newtheorem{remark}[definition]{Remark}

\crefname{definition}{definition}{definitions}
\crefname{axiom}{axiom}{axioms}
\crefname{principle}{principle}{principles}
\crefname{requirement}{requirement}{requirements}
\crefname{proposition}{proposition}{propositions}
\crefname{theorem}{theorem}{theorems}
\crefname{lemma}{lemma}{lemmas}
\crefname{corollary}{corollary}{corollaries}
\crefname{remark}{remark}{remarks}

\newcommand{\kT}{\bm{k}_{T}}
\newcommand{\kiT}{\bm{k}_{iT}}
\newcommand{\pT}{\bm{p}_{T}}
\newcommand{\qT}{\bm{q}_{T}}
\newcommand{\PT}{\bm{P}_{T}}
\newcommand{\DT}{\bm{\Delta}_{T}}
\newcommand{\DNT}{\bm{\Delta}_{N T}}
\newcommand{\bD}{\bm{b}_{\Delta}}
\newcommand{\bM}{\bm{b}_{\mathrm{TMD}}}
\newcommand{\RT}{\bm{R}_{T}}
\newcommand{\dd}{\mathrm{d}}
\newcommand{\Tr}{\operatorname{Tr}}
\newcommand{\tr}{\operatorname{tr}}
\newcommand{\Span}{\operatorname{span}}
\newcommand{\End}{\operatorname{End}}
\newcommand{\Hom}{\operatorname{Hom}}
\newcommand{\Id}{\operatorname{Id}}
\newcommand{\Ker}{\operatorname{Ker}}
\newcommand{\Ran}{\operatorname{Ran}}
\newcommand{\rank}{\operatorname{rank}}
\newcommand{\diag}{\operatorname{diag}}
\newcommand{\Hil}{\mathcal H}
\newcommand{\Phys}{\mathcal H_{\mathrm{phys}}}
\newcommand{\Alg}{\mathcal A}
\newcommand{\Rep}{\mathcal V}
\newcommand{\LF}{\mathrm{LF}}
\newcommand{\TMD}{\mathrm{TMD}}
\newcommand{\GTMD}{\mathrm{GTMD}}
\newcommand{\GPD}{\mathrm{GPD}}
\newcommand{\Amp}{\mathbf{Amp}}
\newcommand{\Dens}{\mathbf{Dens}}
\newcommand{\Match}{\mathbf{Match}}
\newcommand{\Red}{\mathbf{Red}}
\newcommand{\Proc}{\mathbf{Proc}}
\newcommand{\cO}{\mathcal O}
\newcommand{\cM}{\mathcal M}
\newcommand{\cR}{\mathcal R}
\newcommand{\cS}{\mathcal S}
\newcommand{\cV}{\mathcal V}
\newcommand{\cW}{\mathcal W}
\newcommand{\cE}{\mathcal E}
\newcommand{\cG}{\mathcal G}
\newcommand{\cK}{\mathcal K}
\newcommand{\cQ}{\mathcal Q}
\newcommand{\cP}{\mathcal P}
\newcommand{\cC}{\mathcal C}
\newcommand{\cZ}{\mathcal Z}
\newcommand{\cD}{\mathcal D}
\newcommand{\cN}{\mathcal N}
\newcommand{\cF}{\mathcal F}
\newcommand{\cA}{\mathcal A}
\newcommand{\cB}{\mathcal B}
\newcommand{\cU}{\mathcal U}
\newcommand{\eps}{\varepsilon}
\newcommand{\etaW}{\eta_{\!W}}
\newcommand{\bbone}{\bm 1}
\newcommand{\abs}[1]{\left|#1\right|}
\newcommand{\norm}[1]{\left\lVert#1\right\rVert}
\newcommand{\inner}[2]{\left\langle #1\middle|#2\right\rangle}
\newcommand{\avg}[1]{\left\langle #1\right\rangle}
\newcommand{\st}[1]{#1_{\mathrm{ST}}}
\newcommand{\crossT}[2]{\bigl(#1\!\times\!#2\bigr)_z}
\newcommand{\slashn}{\slashed{n}}
\newcommand{\slashnb}{\slashed{\bar n}}

\newcolumntype{Y}{>{\raggedright\arraybackslash}X}
\newcolumntype{L}[1]{>{\raggedright\arraybackslash}p{#1}}

\title{\vspace{-1.0cm}
{\color{navy}\bfseries Volume II: Common Nucleon GTMD Operators\\[2mm]
and Fock-Space Overlaps}\\[4mm]
\large Quark, antiquark, and gluon correlators; zero-skewness recoil;\\
typed reductions; and common-parent validation}
\author{DeuteronWigner project}
\date{28 July 2026}

\begin{document}
\maketitle

\begin{center}
\begin{minipage}{0.94\textwidth}
\colorbox{softblue}{\parbox{0.96\textwidth}{
\textbf{Status and purpose.}
This is the second dynamical volume of the next-level DeuteronWigner
formalism.  Volume~I constructed regulated light-front state spaces, their
physical gauge-constrained subspaces, the block mass operator,
renormalization datum, and controlled truncation system.  The present volume
defines the common nonzero-transfer quark, antiquark, and gluon operator
matrix elements built from those states.  Its primary calculation is the
leading-twist, zero-skewness, positive-
\(x\) DGLAP region for a spin-\(\tfrac12\) nucleon.  It derives an exact
symmetric-frame recoil map, a typed Fock-space overlap kernel, helicity and
transverse-rank projectors, and matched reductions to TMD, GPD, PDF,
form-factor, Wigner, and orbital-angular-momentum objects.  Dynamical
rescattering phases are deferred to Volume~III, the spin-1 nuclear matching
to Volume~IV, and complete perturbative matching/evolution to Volume~V.
}}
\end{minipage}
\end{center}

\begin{abstract}
Generalized transverse-momentum-dependent correlators provide the common
off-forward parent of transverse-momentum-dependent distributions,
generalized parton distributions, collinear densities, form-factor moments,
and partonic Wigner distributions \cite{Meissner2009,BelitskyJiYuan2004,LorcePasquini2011}.  Their use becomes genuinely predictive
only when the quark, antiquark, and gluon correlators are evaluated from one
regulated light-front state rather than supplied as independently normalized
functions.  This volume gives the formal operator and overlap construction
needed for that transition.

We fix a symmetric zero-skewness frame and derive the intrinsic recoil map
for every active and spectator parton.  For an active constituent of fraction
\(x_j\), its incoming and outgoing intrinsic transverse momenta are
\(\bm{k}_{jT}\mp(1-x_j)\DT/2\), whereas a spectator of fraction \(x_i\) carries
\(\bm{k}_{iT}\pm x_i\DT/2\).  The map preserves the intrinsic constraint, has unit
Jacobian, reduces to the identity at zero transfer, and is exchanged by
Hermitian conjugation.  Quark, antiquark, and gluon GTMDs are then formulated
as decorated operator fibers between distinct hadron-momentum Hilbert spaces.
The canonical positive-\(x\), zero-skewness overlap is diagonal in retained
Fock number at zeroth Wilson-line order \cite{Diehl2000}; number-changing contributions are
admitted only through named operator, Wilson-line, zero-mode, effective, or
nonzero-skewness mechanisms.

The leading quark correlator is resolved into vector, axial, and chiral-odd
operator sectors and a complete sixteen-function twist-two GTMD basis.  The
gluon correlator retains an ordered Wilson-link pair, transverse polarization
tensor, and explicit color-flow status.  Antiquarks are represented as
positive-\(x\) active constituents in sea-containing sectors, with
negative-\(x\) relations implemented only through an explicit
charge-conjugation adapter.  Helicity amplitudes and irreducible transverse
harmonics provide convention-safe projectors to named nucleon TMDs.  We
separate the regulated light-front overlap from the ultraviolet- and
rapidity-renormalized, soft-subtracted QCD operator; consequently, the GPD,
PDF, and local-current routes are expressed as matched commuting diagrams,
not as an unconditional integral of a staple-subtracted TMD.  The volume ends
with formal propositions, analytic benchmark models, software interfaces,
property-based tests, and acceptance criteria for passing a common nucleon
parent to the Wilson-line and nuclear volumes.
\end{abstract}

\tableofcontents

\section{Role, scope, and logical status}
\label{sec:scope}

\subsection{The required change in source of prediction}

The accepted pre-evolution model already enforces a common correlator
contract, retains flavor and constituent provenance, and obtains named TMDs
through declared projections.  Its scalar amplitudes, however, still combine
several fits and controlled model sectors.  The next model class must reverse
that direction of inference:
\begin{equation}
 \left\{\cM_{\LF,r}^2,\,\theta_r\right\}
 \longrightarrow
 \left\{|N,P,\Lambda;r\rangle,\,\cO_{a}^{(r)}\right\}
 \longrightarrow
 \left\{W_{q/N}^{(r)},W_{\bar q/N}^{(r)},W_{g/N}^{(r)}\right\}
 \longrightarrow
 \left\{F_A,\cO_{\rm process}\right\}.
\label{eq:source-reversal}
\end{equation}
Here \(r\) is the complete regulator, basis, Fock-sector, and
renormalization label defined in Volume~I.  External PDF, TMD, GPD, lattice,
and form-factor information calibrates or validates the shared state and
operators; it does not provide unrelated central functions.

\begin{principle}[Common-state generation]
\label{prin:common-state}
For a fixed microscopic member and truncation label, every quark flavor,
antiquark flavor, and gluon correlator must be evaluated from the same
normalized eigenstate and compatible operator family.  Removing a fitted
named-TMD table must not make the microscopic calculation undefined.
\end{principle}

\subsection{Primary kinematic domain}

The primary domain of this volume is
\begin{equation}
 \xi=0,
 \qquad
 0<x<1,
 \qquad
 t=-\DT^2\le0,
 \qquad
 \text{leading twist},
\label{eq:primary-domain}
\end{equation}
for a spin-\(\tfrac12\) nucleon.  This domain contains the TMD limit,
transverse Wigner transform, zero-skewness GPDs, and electromagnetic and
energy--momentum form-factor moments.  It is also the cleanest setting for
validating one common recoil convention.

The following are explicit extensions rather than hidden assumptions:
\begin{itemize}
\item \(\xi\ne0\), including ERBL number-changing overlaps and polynomiality;
\item complete Wilson-line rescattering and naive-\(T\)-odd phases;
\item the spin-1 deuteron state and matched nuclear amplitudes;
\item complete soft subtraction, perturbative matching, and multi-scale
      evolution;
\item process factorization and fixed-order \(Y\) terms.
\end{itemize}

\begin{requirement}[No scope migration]
\label{req:no-scope-migration}
A term outside \cref{eq:primary-domain} may enter only with an explicit
mechanism, power counting, operator identity, and replacement rule.  It may
not be absorbed into a fitted scalar GTMD coefficient while retaining the
status of a first-principles overlap.
\end{requirement}

\subsection{Interface with the other volumes}

\begin{longtable}{L{0.16\textwidth}L{0.35\textwidth}L{0.39\textwidth}}
\toprule
Volume & Supplies to Volume II & Receives from Volume II\\
\midrule
\endhead
0 & Typed architecture and fail-closed composition rules & Concrete operator, overlap, reduction, and validation objects.\\
I & Regulated physical Hilbert spaces, normalized Fock states, mass operator, current and truncation data & State-dependent overlap requirements and operator counterterm obligations.\\
III & --- & Bare link-path operators, intermediate-state kernels, and the zero-rescattering parent on which dynamical phases act.\\
IV & --- & Flavor-, helicity-, transfer-, rank-, and link-resolved nucleon GTMD matrices for nuclear matching.\\
V & --- & Regulated GTMD/TMD boundary objects and their operator identities for soft subtraction, matching, and evolution.\\
VI & --- & Shared microscopic sensitivities, common-parent closure residuals, and withheld nucleon observables.\\
\bottomrule
\end{longtable}

\section{Inherited light-front conventions and state normalization}
\label{sec:inherited}

\subsection{Spacetime convention}

We inherit without modification
\begin{align}
 v^\pm&=\frac{v^0\pm v^3}{\sqrt2},
 &v\cdot w&=v^+w^-+v^-w^+-\bm v_T\cdot\bm w_T,
\label{eq:lf-convention}\\
 v^2&=2v^+v^- -\bm v_T^2,
 &\epsilon_T^{12}&=+1.
\label{eq:lf-invariant}
\end{align}
The nucleon states obey the covariant light-front normalization
\begin{equation}
 \langle P',\Lambda';r|P,\Lambda;r\rangle
 =2P^+(2\pi)^3\delta(P'^+-P^+)
 \delta^{(2)}(\bm P'_T-\bm P_T)\delta_{\Lambda'\Lambda}.
\label{eq:hadron-normalization}
\end{equation}

\subsection{Intrinsic Fock variables}

For an \(n\)-parton Fock component,
\begin{equation}
 p_i^+=x_iP^+,
 \qquad
 \bm p_{iT}=x_i\bm P_T+\kiT,
 \qquad
 \sum_{i=1}^n x_i=1,
 \qquad
 \sum_{i=1}^n\kiT=0.
\label{eq:intrinsic-definition}
\end{equation}
The integration measure is
\begin{equation}
 [\dd\Gamma_n]
 =\left[\prod_{i=1}^n
 \frac{\dd x_i\,\dd^2\kiT}{16\pi^3}\right]
 16\pi^3\,
 \delta\!\left(1-\sum_i x_i\right)
 \delta^{(2)}\!\left(\sum_i\kiT\right).
\label{eq:fock-measure}
\end{equation}
A normalized nucleon state has the expansion
\begin{align}
 |N,P,\Lambda;r\rangle
 &=\sum_{\nu\in\cF_r}
 \sum_{\{\alpha_i\}}
 \int[\dd\Gamma_{n_\nu}]
 \frac{\psi^{\Lambda,(r)}_{\nu}
       (\{x_i,\kiT,\alpha_i\})}
      {\sqrt{\prod_i x_i}}
 |\nu;\{x_iP^+,x_i\PT+\kiT,\alpha_i\}\rangle,
\label{eq:nucleon-fock-expansion}\\
 1&=\sum_{\nu\in\cF_r}\sum_{\{\alpha_i\}}
 \int[\dd\Gamma_{n_\nu}]
 \left|\psi^{\Lambda,(r)}_\nu
 (\{x_i,\kiT,\alpha_i\})\right|^2.
\label{eq:fock-normalization}
\end{align}
The composite label \(\alpha_i\) contains species, flavor, color, helicity,
and any basis multiplicity or orbital label required by the regulator.

\begin{definition}[Microscopic member]
\label{def:microscopic-member}
A microscopic member is the tuple
\begin{equation}
 \mathfrak m=
 \left(r,\theta_r,\cM^2_{\LF,r},
 |N;r\rangle,\{\cO_a^{(r)}\},\{J_r^\mu\}\right),
\label{eq:microscopic-member}
\end{equation}
including the state phase convention, counterterms, operator
renormalization data, and numerical convergence manifest.  Member identity
is preserved through every overlap and reduction.
\end{definition}

\begin{requirement}[One normalization ledger]
\label{req:one-normalization}
Quark, antiquark, and gluon overlaps use the same state normalization,
intrinsic measure, sector probabilities, and phase convention.  A species-
specific renormalization of the state is forbidden.  Operator
renormalization and matching factors remain explicit operator data.
\end{requirement}

\section{Symmetric zero-skewness kinematics}
\label{sec:kinematics}

\subsection{Hadron momentum fibers}

Let
\begin{equation}
 P=\frac{p+p'}{2},
 \qquad
 \Delta=p'-p,
 \qquad
 \xi=-\frac{\Delta^+}{2P^+}.
\label{eq:average-transfer}
\end{equation}
At zero skewness choose the symmetric transverse frame
\begin{align}
 \xi&=0,
 &p^+&=p'^+=P^+,
 &\bm P_T&=0,
\label{eq:symmetric-frame-1}\\
 \bm p_T&=-\frac{\DT}{2},
 &\bm p'_T&=+\frac{\DT}{2},
 &t&=\Delta^2=-\DT^2.
\label{eq:symmetric-frame-2}
\end{align}
The on-shell minus components are fixed by
\begin{equation}
 p^- = p'^-=
\frac{M_N^2+\DT^2/4}{2P^+}.
\label{eq:on-shell-minus}
\end{equation}

\begin{definition}[Momentum fiber]
\label{def:momentum-fiber}
For fixed regulator \(r\), the physical nucleon space with total momentum
\(p\) is a fiber \(\Hil_{\mathrm{phys}}(p;r)\).  An off-forward operator is therefore a
map
\begin{equation}
 \cO(P',P):\Hil_{\mathrm{phys}}(p;r)\longrightarrow\Hil_{\mathrm{phys}}(p';r),
\label{eq:operator-fiber-map}
\end{equation}
not a nominal endomorphism of one fixed-momentum space.  A direct-integral
realization may collect the fibers, but their momentum labels remain part of
operator identity.
\end{definition}

\subsection{Average active momentum}

For an active parton \(j\), define its average physical momentum by
\begin{equation}
 \bar k_j^+=x_jP^+,
 \qquad
 \overline{\bm k}_{jT}=\kT.
\label{eq:active-average}
\end{equation}
At zeroth Wilson-line order, the bilocal operator transfers the full hadron
momentum \(\Delta\) through the active line, whereas physical spectator
momenta are unchanged.  Thus
\begin{align}
 \bm p_{jT}^{\rm in}&=\kT-\frac{\DT}{2},
 &\bm p_{jT}^{\rm out}&=\kT+\frac{\DT}{2},
\label{eq:active-physical}\\
 \bm p_{iT}^{\rm in}&=\bm p_{iT}^{\rm out}
 =\kiT,
 &&i\ne j,
\label{eq:spectator-physical}
\end{align}
where the spectator variables on the right are average-frame physical
momenta and satisfy
\(\sum_{i\ne j}\kiT=-\kT\).

\section{Exact zero-skewness recoil map}
\label{sec:recoil}

\subsection{Intrinsic initial and final coordinates}

Applying \cref{eq:intrinsic-definition} separately to the incoming and
outgoing hadrons gives the following map.

\begin{definition}[Symmetric-frame recoil map]
\label{def:recoil-map}
For active parton \(j\),
\begin{align}
 \bm\kappa_{jT}^{\rm in}
 &=\kT-\frac{1-x_j}{2}\DT,
 &\bm\kappa_{jT}^{\rm out}
 &=\kT+\frac{1-x_j}{2}\DT,
\label{eq:active-recoil-map}\\
 \bm\kappa_{iT}^{\rm in}
 &=\kiT+\frac{x_i}{2}\DT,
 &\bm\kappa_{iT}^{\rm out}
 &=\kiT-\frac{x_i}{2}\DT,
 \qquad i\ne j.
\label{eq:spectator-recoil-map}
\end{align}
The longitudinal fractions are unchanged:
\begin{equation}
 x_i^{\rm in}=x_i^{\rm out}=x_i.
\label{eq:longitudinal-recoil}
\end{equation}
We denote the full affine transformation by
\(\cR_{j,\Delta}^{(n)}\).
\end{definition}

\begin{proposition}[Intrinsic closure]
\label{prop:intrinsic-closure}
The map in \cref{eq:active-recoil-map,eq:spectator-recoil-map} preserves the
intrinsic transverse constraints:
\begin{equation}
 \sum_i\bm\kappa_{iT}^{\rm in}=0,
 \qquad
 \sum_i\bm\kappa_{iT}^{\rm out}=0.
\label{eq:recoil-closure}
\end{equation}
\end{proposition}

\begin{proof}
Using \(\sum_i\kiT=0\) and
\(\sum_{i\ne j}x_i=1-x_j\), the incoming transfer contribution is
\begin{equation}
 -\frac{1-x_j}{2}\DT
 +\sum_{i\ne j}\frac{x_i}{2}\DT=0.
\end{equation}
The outgoing result has the opposite signs and vanishes identically.
\end{proof}

\begin{proposition}[Physical momentum assignment]
\label{prop:physical-assignment}
When the intrinsic variables of \cref{def:recoil-map} are recombined with the
incoming and outgoing hadron momenta, the active physical parton receives
\(\DT\) and every spectator physical momentum is unchanged.
\end{proposition}

\begin{proof}
For the active line,
\begin{align}
 x_j\left(-\frac{\DT}{2}\right)
 +\bm\kappa_{jT}^{\rm in}
 &=\kT-\frac{\DT}{2},\\
 x_j\left(+\frac{\DT}{2}\right)
 +\bm\kappa_{jT}^{\rm out}
 &=\kT+\frac{\DT}{2}.
\end{align}
For a spectator,
\begin{align}
 x_i\left(-\frac{\DT}{2}\right)
 +\bm\kappa_{iT}^{\rm in}&=\kiT,\\
 x_i\left(+\frac{\DT}{2}\right)
 +\bm\kappa_{iT}^{\rm out}&=\kiT.
\end{align}
\end{proof}

\begin{proposition}[Unit Jacobian and involution]
\label{prop:recoil-jacobian}
At fixed \(x_i\) and \(\DT\), the recoil map is affine with unit Jacobian.
Moreover,
\begin{equation}
 \cR_{j,-\Delta}^{(n)}
 \left(\{\bm\kappa_i^{\rm out}\},
       \{\bm\kappa_i^{\rm in}\}\right)
 =
 \left(\{\bm\kappa_i^{\rm in}\},
       \{\bm\kappa_i^{\rm out}\}\right),
\label{eq:recoil-involution}
\end{equation}
and at \(\DT=0\) both coordinate sets reduce to the same intrinsic point.
\end{proposition}

\begin{proof}
The linear part of each translation is the identity.  Reversing \(\DT\)
interchanges every incoming and outgoing shift.  The forward statement is
immediate.
\end{proof}

\begin{requirement}[Single recoil authority]
\label{req:single-recoil-authority}
Every quark, antiquark, and gluon overlap in this volume calls the same
versioned recoil-map object.  A projector, observable, or nuclear adapter may
not reproduce the shifts independently.  Alternative frames require a new
named recoil convention and cross-convention tests.
\end{requirement}

\subsection{Typed recoil identity}

The computational identity is
\begin{equation}
 \mathfrak R_{j,\Delta}=
 \left(
 \texttt{SYMMETRIC\_XI0},
 n,j,\{x_i\},\DT,
 \text{incoming fiber},
 \text{outgoing fiber},
 \text{Jacobian}=1
 \right).
\label{eq:typed-recoil-id}
\end{equation}
The implementation must distinguish the average active momentum \(\kT\),
the intrinsic spectator coordinates \(\kiT\), and the total transfer \(\DT\).
They have equal physical dimensions but different semantic types.

\section{Decorated operator identity and map classes}
\label{sec:operator-identity}

\subsection{Operator identity tuple}

\begin{definition}[Decorated GTMD operator]
\label{def:decorated-operator}
A production operator is identified by
\begin{equation}
 \mathfrak o=
 \Bigl(
 a,f,\Gamma,
 p,p',\xi,t,
 \gamma, R_{\rm color},
 \mathfrak l_g,\mathfrak c_g,
 \cR_{\rm UV},\cR_{\rm rap},
 \cS_{\rm soft},\mu,\zeta,
 \mathfrak s,\mathfrak r,
 r
 \Bigr).
\label{eq:decorated-op-tuple}
\end{equation}
Here \(a\in\{q,\bar q,g\}\), \(f\) is the explicit flavor where
applicable, \(\Gamma\) is the Dirac or transverse Lorentz projection,
\(\gamma\) is the geometric Wilson path, \(R_{\rm color}\) its
representation, \(\mathfrak l_g\) an ordered gluon-link pair,
\(\mathfrak c_g\) a gluonic color-contraction status, \(\mathfrak s\) the
operator scheme, \(\mathfrak r\) the transverse-rank convention, and \(r\)
the microscopic truncation label.
\end{definition}

The color-contraction status takes values
\begin{equation}
 \mathfrak c_g\in
 \{\texttt{NOT\_APPLICABLE},
   \texttt{UNRESOLVED},
   \texttt{F\_TYPE},
   \texttt{D\_TYPE},
   \texttt{PROCESS\_WEIGHTED}\}.
\label{eq:gluon-color-status}
\end{equation}
It is incorrect to label every generic T-even gluon GTMD as either
\(f\)- or \(d\)-type.  Those independent contractions become indispensable
for gluonic-pole and naive-\(T\)-odd structures and for process-specific link
combinations.  The status field prevents both premature identification and
silent omission.

\begin{requirement}[Fail-closed operator use]
\label{req:fail-closed-operator}
An operator with an \texttt{UNRESOLVED} field may be stored and inspected,
but it may not enter a calculation that mathematically requires that field.
For example, a T-even zero-rescattering benchmark may leave
\(\mathfrak c_g\) unresolved; a gluon Sivers process projection may not.
\end{requirement}

\subsection{Bare, regulated, and renormalized objects}

We distinguish three layers:
\begin{align}
 \cW_{\rm bare}[\gamma]
 &:\text{formal gauge-covariant bilocal operator},
\label{eq:bare-operator}\\
 W_{\LF}^{(r)}
 &=\langle N;r|\cW_{\rm reg}^{(r)}|N;r\rangle,
\label{eq:regulated-overlap}\\
 W_{\rm QCD}(\mu,\zeta;\mathfrak s)
 &=\cZ_{\LF\to\rm QCD}^{(r\to\mathfrak s)}
 \otimes W_{\LF}^{(r)}
 +\delta W_{\rm trunc}^{(r)}.
\label{eq:matched-qcd-gtmd}
\end{align}
For a staple-like GTMD/TMD definition, \(\cZ\) includes ultraviolet
renormalization, rapidity renormalization, and soft subtraction
\cite{Collins2011,Echevarria2016}.  It may also
mix operator structures.  The low-resolution wave-function overlap is
therefore not itself a TMD in a declared QCD scheme.

\begin{principle}[No nominal equality across layers]
\label{prin:no-nominal-equality}
The symbols in \cref{eq:regulated-overlap,eq:matched-qcd-gtmd} may have the
same kinematic arguments, but they are different typed objects.  Equality is
replaced by an explicit matching morphism with a stated accuracy and
uncertainty.
\end{principle}

\subsection{Map-class separation}

The present calculation uses the same map classes defined in Volume~0:
\begin{longtable}{L{0.15\textwidth}L{0.31\textwidth}L{0.44\textwidth}}
\toprule
Class & Examples in Volume II & Composition rule\\
\midrule
\endhead
\(\Amp\) & Recoil embeddings, active-parton kernels, overlap contractions & Linear amplitude maps compose only between compatible momentum fibers and symmetry blocks.\\
\(\Dens\) & Forward spectator trace and Gram correlator & Positive or completely positive maps arise only after the amplitude is formed and an unresolved subsystem is traced.\\
\(\Match\) & Regulator-to-QCD matching, soft subtraction, link shortening, small-separation OPE & Scheme-changing maps carry accuracy, scale, and mixing metadata.\\
\(\Red\) & TMD limit, Wigner transform, helicity projector, GPD/PDF/current moment & Typed limits, Fourier transforms, integrals, and moments act only on compatible operator spaces.\\
\(\Proc\) & SIDIS/DY/gluon-process link selection and hard contraction & Deferred to process volumes; process weights may not be inserted into the parent operator silently.\\
\bottomrule
\end{longtable}

\begin{requirement}[Explicit adapters]
\label{req:explicit-adapters}
A composition across map classes requires a registered adapter that states
its domain, codomain, assumptions, order, and error estimate.  Array-shape
compatibility is never sufficient.
\end{requirement}

\section{Common Fock-space overlap architecture}
\label{sec:common-overlap}

\subsection{General operator block form}

Let \(P_\nu\) project onto a retained Fock sector.  Any regulated bilocal
operator has blocks
\begin{equation}
 \cO^{(r)}_{\nu\mu}=P_\nu\cO^{(r)}P_\mu.
\label{eq:operator-block}
\end{equation}
Its matrix element is
\begin{equation}
 W^{(r)}_{\Lambda'\Lambda}
 =\sum_{\nu,\mu\in\cF_r}
 \inner{\psi_{\nu,\Lambda'}^{\rm out}}
        {\cK^{(r)}_{\nu\mu}
         \psi_{\mu,\Lambda}^{\rm in}},
\label{eq:general-block-overlap}
\end{equation}
where \(\cK_{\nu\mu}\) contains the active operator, spectator matching,
recoil map, color and helicity intertwiners, Wilson-line insertions, and all
normalization factors inherited from the field expansion.

\begin{definition}[Canonical one-body core]
\label{def:canonical-one-body}
The canonical one-body core is the contribution in the primary domain for
which:
\begin{enumerate}
\item the leading good-field bilocal annihilates and recreates one active
      constituent of the same species;
\item the retained Fock sector is unchanged, \(\nu=\mu\);
\item all spectators are matched in species, flavor, color, helicity, and
      physical momentum;
\item the geometric link is evaluated at zeroth dynamical rescattering order,
      while its path identity remains stored.
\end{enumerate}
\end{definition}

For this core,
\begin{align}
 W_{a/N,\Lambda'\Lambda}^{[J],(r)}
 (x,\kT,\DT)
 &=\sum_{\nu\in\cF_r}
   \sum_{j\in a(\nu)}
   \sum_{\boldsymbol\alpha',\boldsymbol\alpha}
   \int[\dd\Gamma_{n_\nu}]
   \delta(x-x_j)
   \delta^{(2)}(\kT-\bm{k}_{jT})
\nonumber\\
 &\quad\times
 \psi_{\nu,\Lambda'}^{(r)*}
 \bigl(\{x_i,\bm\kappa_{iT}^{\rm out},\alpha_i'\}\bigr)
 \cK_{a,j;\boldsymbol\alpha'\boldsymbol\alpha}^{[J],(r)}
 \psi_{\nu,\Lambda}^{(r)}
 \bigl(\{x_i,\bm\kappa_{iT}^{\rm in},\alpha_i\}\bigr).
\label{eq:master-diagonal-overlap}
\end{align}
The active index set \(a(\nu)\) is species- and flavor-resolved.  The kernel
contains spectator Kronecker deltas and an active spin-color matrix:
\begin{equation}
 \cK_{a,j}^{[J]}
 =\cN_{a,J}^{(r)}
 \left[\prod_{i\ne j}
 \delta_{\alpha_i'\alpha_i}\right]
 \cP_{a,j}^{[J]}
 \cC_{a,j}^{[J]}.
\label{eq:active-kernel-factor}
\end{equation}
The normalization \(\cN_{a,J}^{(r)}\) is fixed by the field normalization and
verified against number, helicity, and momentum moments; it is not an
adjustable TMD normalization.

\begin{proposition}[Forward Gram form]
\label{prop:forward-gram}
At \(\DT=0\), for a number-density projection with a positive active
polarization projector, the canonical one-body correlator is a Gram matrix in
nucleon and active-parton helicity space.
\end{proposition}

\begin{proof}
At zero transfer the incoming and outgoing wave functions are evaluated at
the same intrinsic point.  Diagonalizing the positive active projector gives
\(\cP=R^\dagger R\).  Absorbing \(R\) into the active amplitude and collecting
all spectator labels into \(X\), one obtains
\begin{equation}
 \Phi_{\beta\alpha}(x,\kT)
 =\sum_X A_{\beta X}(x,\kT)A_{\alpha X}^*(x,\kT),
\end{equation}
which is Hermitian and positive semidefinite.
\end{proof}

\subsection{When number-changing blocks are allowed}

The diagonal core is not a claim that every GTMD operator is forever diagonal
in Fock number.  Number-changing blocks arise through named mechanisms:
\begin{longtable}{L{0.22\textwidth}L{0.35\textwidth}L{0.33\textwidth}}
\toprule
Mechanism & Operator/Fock effect & Required declaration\\
\midrule
\endhead
Nonzero skewness ERBL region & Pair emission or annihilation and \(n+1\leftrightarrow n-1\) overlaps & \(\xi\ne0\), support region, exact recoil map, polynomiality tests.\\
Wilson-line expansion & Eikonal gluon insertions and intermediate-state denominators & Link path, perturbative/eikonal order, rapidity regulator, pole prescription, color tensor.\\
Constrained fields and zero modes & Instantaneous or endpoint operators & Gauge fixing, zero-mode prescription, power counting, counterterm partner.\\
Feshbach/effective operator & Excluded sectors reappear as induced kernels & Projection space, matching energy/resolution, subtraction and replacement identity.\\
Renormalization mixing & Operator structures mix under UV or rapidity renormalization & Mixing matrix, scheme, scale, order, truncation error.\\
\bottomrule
\end{longtable}

\begin{requirement}[Named off-diagonal source]
\label{req:named-off-diagonal}
Every \(\nu\ne\mu\) block in \cref{eq:general-block-overlap} must identify
one row of the preceding table or an equally explicit physical mechanism.
An off-diagonal block may not be introduced solely to populate a missing
named TMD.
\end{requirement}

\subsection{Spectator identity and permutation multiplicity}

For identical constituents the sum over active labels and the statistical
normalization of the Fock basis must not double count the same physical
configuration.  Two equivalent implementations are admissible:
\begin{enumerate}
\item explicitly sum over every labeled active constituent using a normalized
      antisymmetric/symmetric basis;
\item select a canonical active slot and multiply by a verified occupation
      number.
\end{enumerate}
They must agree in analytic benchmarks.  The active-selection identity is
\begin{equation}
 \mathfrak a_j=
 (\nu,j,a,f,\lambda_j,c_j,\text{permutation orbit},\text{multiplicity}).
\label{eq:active-selection-id}
\end{equation}

\begin{proposition}[Flavor-number moment]
\label{prop:flavor-number}
For a conserved vector current and a state with no unresolved current
counterterm, integrating the forward positive-\(x\) quark-minus-antiquark
number correlator gives the net flavor charge of the state.
\end{proposition}

\begin{proof}
The \(x\) and \(\kT\) integrals remove the active delta functions.  The sum
over active quarks counts quark occupation; the antiquark operator contributes
with the charge-conjugate sign fixed by the local vector current.  State
normalization then yields the expectation value of the flavor number
operator.  At finite truncation, any violation diagnoses an omitted current
counterterm, incomplete sector normalization, or an inconsistent active
multiplicity.
\end{proof}

\section{Quark GTMD operator and overlap}
\label{sec:quark}

\subsection{Gauge-covariant bilocal operator}

For quark flavor \(q\), define the regulated unsubtracted correlator
\begin{align}
 W_{\Lambda'\Lambda}^{q[\Gamma],(r)}
 (x,\kT,\DT;\gamma)
 &=\frac12
 \int\frac{\dd z^-\,\dd^2\bm z_T}{(2\pi)^3}
 e^{ixP^+z^- -i\kT\cdot\bm z_T}
\nonumber\\
 &\quad\times
 \langle p',\Lambda';r|
 \bar\psi_q(-z/2)\,\Gamma\,
 U_\gamma[-z/2,z/2]\,
 \psi_q(z/2)
 |p,\Lambda;r\rangle
 \big|_{z^+=0}.
\label{eq:quark-gtmd-definition}
\end{align}
At leading twist,
\begin{equation}
 \Gamma\in
 \left\{\gamma^+,\ \gamma^+\gamma_5,
 i\sigma^{j+}\gamma_5\right\},
 \qquad j=1,2.
\label{eq:quark-leading-projectors}
\end{equation}
The path \(\gamma\) is part of operator identity even when the dynamical link
is set to unity in the zero-rescattering pilot.

\subsection{Good-field reduction}

Let
\begin{equation}
 \psi_+=\Lambda_+\psi,
 \qquad
 \Lambda_+=\frac12\gamma^-\gamma^+.
\label{eq:good-field-projector}
\end{equation}
After fixing the light-front spinor phase convention, the normalized active
helicity kernels are represented on the two-component good-spin space by
\begin{align}
 \cP_q^{[\gamma^+]}&=\bbone,
\label{eq:quark-unpol-kernel}\\
 \cP_q^{[\gamma^+\gamma_5]}&=\sigma_3,
\label{eq:quark-helicity-kernel}\\
 \cP_q^{[i\sigma^{j+}\gamma_5]}&=\sigma_T^j,
\label{eq:quark-transverse-kernel}
\end{align}
where \(\sigma_T^1,\sigma_T^2\) carry the declared transverse-spin phase
adapter.  These equations specify the normalized helicity action; all
\(P^+\), \(x_j\), and field-expansion factors remain in
\(\cN_{q,J}^{(r)}\) and are fixed by the local-current limits.

\begin{requirement}[Spinor convention adapter]
\label{req:spinor-adapter}
The code must validate \crefrange{eq:quark-unpol-kernel}{eq:quark-transverse-kernel}
against explicit light-front spinors in the selected gamma-matrix convention.
A Pauli matrix may not be used merely because its array shape is correct.
\end{requirement}

\subsection{Canonical quark overlap}

Substituting the quark active kernel into
\cref{eq:master-diagonal-overlap} gives
\begin{align}
 W_{\Lambda'\Lambda}^{q[\Gamma],(r)}
 (x,\kT,\DT)
 &=\sum_{\nu}
   \sum_{j\in q(\nu)}
   \int[\dd\Gamma_{n_\nu}]
   \delta(x-x_j)\delta^{(2)}(\kT-\bm{k}_{jT})
\nonumber\\
 &\quad\times
 \sum_{\boldsymbol\lambda',\boldsymbol\lambda}
 \sum_{\boldsymbol c',\boldsymbol c}
 \psi_{\nu,\Lambda'}^{(r)*}
 \left(\{x_i,\bm\kappa_{iT}^{\rm out},
             \lambda_i',c_i'\}\right)
\nonumber\\
 &\quad\times
 \left[\prod_{i\ne j}
 \delta_{\lambda_i'\lambda_i}\delta_{c_i'c_i}\right]
 \left[\cP_q^{[\Gamma]}\right]_{\lambda_j'\lambda_j}
 \delta_{c_j'c_j}
 \psi_{\nu,\Lambda}^{(r)}
 \left(\{x_i,\bm\kappa_{iT}^{\rm in},
             \lambda_i,c_i\}\right)
 \cN_{q,\Gamma}^{(r)}.
\label{eq:quark-master-overlap}
\end{align}
Flavor is selected exactly by \(j\in q(\nu)\).  The same equation therefore
produces distinct \(u\), \(d\), and heavier-flavor objects without imposing
an isospin identification.

\begin{theorem}[Regulated common-parent nesting]
\label{thm:canonical-common-parent}
Assume the state and active kernel satisfy the Volume~I normalization and
current conditions, the recoil map is \cref{def:recoil-map}, and every route
retains the same regulated operator identity.  Then the forward limit and the
formal transverse-momentum and longitudinal moments of
\cref{eq:quark-master-overlap} are nested linear reductions of one regulated
matrix element.  A physical GPD, PDF, or local current is obtained only after
the declared link-shortening, renormalization, and matching maps of
\cref{sec:reductions}.  No additional normalization is permitted at the
named-function level.
\end{theorem}

\begin{proof}
The forward limit changes only the recoil map, setting its incoming and
outgoing arguments equal.  The \(\kT\) integral removes the active transverse
delta function without changing the state or active kernel, and the
subsequent \(x\) integral removes the longitudinal delta function.  These are
nested regulated coordinate reductions of the same overlap.  Whether the
resulting regulated bilocal equals a straight-link collinear or local QCD
operator is a separate operator-matching statement, carried by an explicit
map.  Because the integration measure, state normalization, active
multiplicity, and matching identity remain fixed along each declared route,
an extra named-function factor would alter a number, current, or matched
moment and violate the hypotheses.
\end{proof}

\subsection{Scalar GTMD basis}
\label{sec:quark-basis}

At twist two the spin-\(\tfrac12\) quark correlator contains sixteen complex
scalar GTMDs, conventionally grouped as \cite{Meissner2009,LorcePasquini2013}
\begin{equation}
 \{F_{1,n}\}_{n=1}^4,
 \qquad
 \{G_{1,n}\}_{n=1}^4,
 \qquad
 \{H_{1,n}\}_{n=1}^8.
\label{eq:sixteen-quark-gtmds}
\end{equation}
They depend on
\begin{equation}
 \left(x,\xi,\kT^2,\kT\cdot\DT,\DT^2;\gamma,\mu,\zeta,\mathfrak s,r\right).
\label{eq:gtmd-scalar-arguments}
\end{equation}
At \(\xi=0\), define a complete linearly independent operator basis
\begin{align}
 \cB_V&=\{B_{V,1},\ldots,B_{V,4}\},
 &\cB_A&=\{B_{A,1},\ldots,B_{A,4}\},
 &\cB_T^j&=\{B_{T,1}^j,\ldots,B_{T,8}^j\},
\label{eq:quark-bases}
\end{align}
constructed from nucleon spinors, \(\kT\), \(\DT\),
\(g_T^{ij}\), \(\epsilon_T^{ij}\), and the leading-twist Dirac bilinears.
Then
\begin{align}
 W^{[\gamma^+]}&=\sum_{n=1}^4F_{1,n}B_{V,n},
\label{eq:vector-basis-expansion}\\
 W^{[\gamma^+\gamma_5]}&=\sum_{n=1}^4G_{1,n}B_{A,n},
\label{eq:axial-basis-expansion}\\
 W^{[i\sigma^{j+}\gamma_5]}&=\sum_{n=1}^8H_{1,n}B_{T,n}^j.
\label{eq:tensor-basis-expansion}
\end{align}

A concrete Meissner--Metz--Schlegel convention adapter may be installed, but
the production projector library is generated from the basis Gram matrix,
not from hand-coded inversions.  For a fixed non-degenerate kinematic point,
\begin{equation}
 G_{AB}=\sum_{\Lambda'\Lambda,J}
 \left(B_A^{[J]}\right)_{\Lambda'\Lambda}^*
 \left(B_B^{[J]}\right)_{\Lambda'\Lambda},
\label{eq:gtmd-gram}
\end{equation}
and
\begin{equation}
 \cP_A=\sum_B(G^{-1})_{AB}B_B,
 \qquad
 F_A=\langle\cP_A,W\rangle.
\label{eq:gtmd-projector}
\end{equation}
At degenerate points such as \(\DT=0\), the rank drops because distinct GTMDs
collapse onto fewer TMD combinations.  The implementation must project
before taking the limit or use the correct reduced basis.

\begin{requirement}[Basis provenance]
\label{req:basis-provenance}
Every scalar GTMD carries the reference basis, gamma-matrix convention,
spinor phase, mass factors, Fourier phase, and any linear transformation to a
published convention.  The name \(F_{1,4}\), for example, is not sufficient
operator identity.
\end{requirement}

\subsection{Helicity-amplitude representation}

Define active-quark helicity amplitudes
\begin{equation}
 H^q_{\Lambda'\lambda',\Lambda\lambda}
 =\frac12\left[
 W^{q[\gamma^+]}_{\Lambda'\Lambda}
 +\lambda W^{q[\gamma^+\gamma_5]}_{\Lambda'\Lambda}
 \right]\delta_{\lambda'\lambda}
 +H^q_{\rm flip},
\label{eq:quark-helicity-amplitudes}
\end{equation}
where \(H^q_{\rm flip}\) is obtained from
\(i\sigma^{j+}\gamma_5\).  The helicity representation is primary for
computational projection because it connects directly to the Fock amplitudes
and makes forward positivity transparent.

For the helicity-conserving diagonal amplitudes, a useful convention adapter
has the schematic structure
\begin{equation}
 H_{\Lambda\lambda,\Lambda\lambda}
 =\frac12\left[
 F_{1,1}+\Lambda\lambda G_{1,4}
 +\frac{i\,\crossT{\kT}{\DT}}{M_N^2}
 \left(\Lambda F_{1,4}-\lambda G_{1,1}\right)
 \right],
\label{eq:f14-g11-helicity}
\end{equation}
with \(\Lambda,\lambda=\pm1\) as signed helicity labels and with phases fixed
by the chosen standard adapter.  This equation exhibits why \(F_{1,4}\) and
\(G_{1,1}\) require nonzero transfer and orbital interference.  It is a
basis-adapter relation, not a license to infer either function from a TMD
alone.

\begin{remark}[No model-independent GPD--TMD shortcut]
Sharing a common GTMD ancestor does not imply an algebraic relation between an
arbitrary GPD and TMD.  Different reductions can select different scalar
GTMD combinations, and Wilson-line/matching data can differ.  The valid
statement is common ancestry plus explicit reduction, not a universal
lensing identity.
\end{remark}

\section{Antiquarks as explicit positive-x constituents}
\label{sec:antiquark}

\subsection{Primary representation}

Antiquarks are stored as separate positive-\(x\) objects,
\begin{equation}
 W_{\bar q/N}^{[\Gamma]}(x,\kT,\DT),
 \qquad 0<x<1,
\label{eq:positive-x-antiquark}
\end{equation}
and are evaluated by \cref{eq:master-diagonal-overlap} with the active sum
restricted to explicit antiquarks:
\begin{equation}
 j\in\bar q(\nu).
\label{eq:active-antiquark-set}
\end{equation}
Consequently, a valence-only \(|qqq\rangle\) state produces no antiquark
parent.  The minimum diagonal support is supplied by sectors such as
\begin{equation}
 |qqqq\bar q\rangle,
 \qquad
 |qqqq\bar q g\rangle,
 \qquad\ldots.
\label{eq:antiquark-sectors}
\end{equation}
Sea distributions generated only by an effective operator or an integrated-
out sector must be tagged as induced contributions rather than diagonal
antiquark overlaps.

\subsection{Charge-conjugation adapter}

The quark bilocal over \(-1<x<1\) and the separate positive-\(x\) quark and
antiquark objects are two representations of the same underlying quark field
operator.  Their conversion is controlled by
\begin{equation}
 C\Gamma^T C^{-1}=\chi_\Gamma\Gamma,
 \qquad
 \chi_\Gamma\in\{+1,-1\},
\label{eq:charge-conjugation-gamma}
\end{equation}
together with Wilson-path reversal, endpoint exchange, color conjugation,
and the selected Fourier convention.  We therefore define a typed adapter
\begin{equation}
 \cA_C^{[\Gamma]}:
 W_q^{[\Gamma]}(x<0;\gamma)
 \longrightarrow
 W_{\bar q}^{[\Gamma]}(-x;\gamma_C),
\label{eq:charge-conjugation-adapter}
\end{equation}
whose sign and phase are derived from \cref{eq:charge-conjugation-gamma}.

\begin{requirement}[No silent negative-\(x\) identity]
\label{req:no-silent-negative-x}
The production representation is \cref{eq:positive-x-antiquark}.  A
negative-\(x\) relation may be used only through
\(\cA_C^{[\Gamma]}\), with separate tests for vector, axial, and chiral-odd
operators.  One sign may not be copied across all Dirac structures.
\end{requirement}

\subsection{Quark--antiquark ledgers}

For each flavor, define the regulated forward moments
\begin{align}
 N_q^{(r)}&=\int_0^1\dd x\int\dd^2\kT\,
 \left[f_1^{q,(r)}(x,\kT)-f_1^{\bar q,(r)}(x,\kT)\right],
\label{eq:net-flavor-number}\\
 \Sigma_q^{(r)}&=\int_0^1\dd x\int\dd^2\kT\,
 \left[f_1^{q,(r)}(x,\kT)+f_1^{\bar q,(r)}(x,\kT)\right],
\label{eq:total-flavor-occupation}\\
 \Delta q^{(r)}&=\int_0^1\dd x\int\dd^2\kT\,
 \left[g_{1L}^{q,(r)}(x,\kT)+g_{1L}^{\bar q,(r)}(x,\kT)\right].
\label{eq:flavor-helicity-ledger}
\end{align}
Only \(N_q\) is a conserved flavor-number charge.  The occupation and
helicity moments depend on resolution, operator scheme, and matching.  Their
truncation behavior is nevertheless a required diagnostic.

\begin{proposition}[Sea-sector zero test]
\label{prop:sea-zero-test}
If all retained Fock sectors contain no explicit antiquark of flavor \(q\)
and the antiquark effective-operator blocks are disabled, then the canonical
positive-\(x\) antiquark overlap vanishes identically.
\end{proposition}

\begin{proof}
The active set \(\bar q(\nu)\) is empty in every term of
\cref{eq:master-diagonal-overlap}; hence the active sum is empty.
\end{proof}

\section{Gluon GTMD operator and overlap}
\label{sec:gluon}

\subsection{Ordered-link field-strength correlator}

For the gluon parent, define \cite{MuldersRodrigues2001,LorcePasquini2013,Boer2016}
\begin{align}
 W_{\Lambda'\Lambda}^{g,ij,(r)}
 (x,\kT,\DT;\gamma_1,\gamma_2)
 &=\frac{1}{xP^+}
 \int\frac{\dd z^-\,\dd^2\bm z_T}{(2\pi)^3}
 e^{ixP^+z^- -i\kT\cdot\bm z_T}
\nonumber\\
 &\quad\times
 \langle p',\Lambda';r|
 \Tr_c\!\Big[
 F^{+i}(-z/2)U_{\gamma_1}[-z/2,z/2]
\nonumber\\[-1mm]
 &\hspace{4.0cm}\times
 F^{+j}(z/2)U_{\gamma_2}[z/2,-z/2]
 \Big]
 |p,\Lambda;r\rangle
 \big|_{z^+=0}.
\label{eq:gluon-gtmd-definition}
\end{align}
The prefactor is part of the declared normalization convention.  Alternative
normalizations are admissible only through an explicit linear adapter and
must reproduce the same collinear gluon momentum fraction.

\begin{definition}[Ordered gluon-link identity]
\label{def:ordered-gluon-link}
The gluon link object is
\begin{equation}
 \mathfrak l_g=([\gamma_1],[\gamma_2],R_A,
 \text{endpoint order},\text{staple directions},
 \text{transverse closures},\text{cusp data}).
\label{eq:ordered-gluon-link}
\end{equation}
The pair is ordered.  Interchanging its entries is not a harmless display
operation and may change charge-conjugation and process properties.
\end{definition}

\begin{requirement}[No quark-link analogy]
\label{req:no-quark-link-analogy}
A gluon link pair, color contraction, or process sign may not be inferred from
a quark staple by name matching.  The process object must select the ordered
pair and any \(f^{abc}\)- or \(d^{abc}\)-type gluonic-pole combination
explicitly.
\end{requirement}

\subsection{Transverse polarization decomposition}

In two-dimensional Euclidean transverse notation, decompose any gluon matrix
as
\begin{equation}
 W_g^{ij}
 =\frac12\delta_T^{ij}\,W_U
 -\frac{i}{2}\epsilon_T^{ij}\,W_H
 +W_L^{ij},
 \qquad
 W_L^{ij}=W_L^{ji},
 \qquad
 \delta_{T,ij}W_L^{ij}=0.
\label{eq:gluon-polarization-decomp}
\end{equation}
The corresponding projectors are
\begin{align}
 W_U&=\delta_{T,ij}W_g^{ij},
\label{eq:gluon-unpol-projector}\\
 W_H&=i\epsilon_{T,ij}W_g^{ij},
\label{eq:gluon-helicity-projector}\\
 W_L^{ij}&=\left[
 \frac12(\delta^i_k\delta^j_l+\delta^i_l\delta^j_k)
 -\frac12\delta_T^{ij}\delta_{T,kl}
 \right]W_g^{kl}.
\label{eq:gluon-linear-projector}
\end{align}
The scalar normalization is anchored in the forward limit by a declared
convention such as
\begin{equation}
 \Gamma_g^{ij}(x,\kT)
 =\frac{x}{2}\left[
 \delta_T^{ij} f_1^g
 +\frac{k_T^{ij}}{M_N^2}h_1^{\perp g}
 +\cdots\right],
 \qquad
 k_T^{ij}=k_T^ik_T^j-\frac12\delta_T^{ij}\kT^2.
\label{eq:gluon-forward-normalization}
\end{equation}
Then
\begin{equation}
 f_1^g(x,\kT)=\frac1x\delta_{T,ij}\Gamma_g^{ij}(x,\kT)
\label{eq:gluon-f1-extraction}
\end{equation}
and
\begin{equation}
 \int_0^1\dd x\,x
 \left[\sum_q(q+\bar q)(x;\mu)+g(x;\mu)\right]=1
\label{eq:momentum-sum-rule}
\end{equation}
after matching to a common collinear scheme.

\subsection{Canonical gluon overlap}

The active sum runs over explicit gluons in sectors such as
\begin{equation}
 |qqqg\rangle,
 \quad |qqqgg\rangle,
 \quad |qqqq\bar q g\rangle,
 \quad\ldots.
\label{eq:gluon-containing-sectors}
\end{equation}
At zeroth rescattering order,
\begin{align}
 W_{\Lambda'\Lambda}^{g,ij,(r)}
 (x,\kT,\DT)
 &=\sum_{\nu}
   \sum_{j\in g(\nu)}
   \int[\dd\Gamma_{n_\nu}]
   \delta(x-x_j)\delta^{(2)}(\kT-\bm{k}_{jT})
\nonumber\\
 &\quad\times
 \sum_{\boldsymbol\lambda',\boldsymbol\lambda}
 \sum_{\boldsymbol c',\boldsymbol c}
 \psi_{\nu,\Lambda'}^{(r)*}
 \left(\{x_i,\bm\kappa_{iT}^{\rm out},
             \lambda_i',c_i'\}\right)
\nonumber\\
 &\quad\times
 \left[\prod_{\ell\ne j}
 \delta_{\lambda_\ell'\lambda_\ell}
 \delta_{c_\ell'c_\ell}\right]
 \left[\cP_g^{ij}\right]_{\lambda_j'\lambda_j}
 \delta_{c_j'c_j}
 \psi_{\nu,\Lambda}^{(r)}
 \left(\{x_i,\bm\kappa_{iT}^{\rm in},
             \lambda_i,c_i\}\right)
 \cN_{g}^{(r)}.
\label{eq:gluon-master-overlap}
\end{align}
On the physical transverse-helicity space, the active polarization kernel can
be constructed from
\begin{equation}
 \epsilon_{\lambda'}^{i*}\epsilon_\lambda^j
\label{eq:gluon-active-kernel}
\end{equation}
and projected with
\crefrange{eq:gluon-unpol-projector}{eq:gluon-linear-projector}.  The normalization
\(\cN_g^{(r)}\) is fixed by the field-strength convention and the momentum
sum rule, not by comparison with a named gluon TMD curve.

\begin{proposition}[Explicit-gluon zero test]
\label{prop:gluon-zero-test}
If no retained sector contains an explicit gluon and no induced gluon
operator is enabled, then the canonical diagonal gluon overlap vanishes.
\end{proposition}

\begin{proof}
For every retained sector, the active set \(g(\nu)\) in
\cref{eq:gluon-master-overlap} is empty.
\end{proof}

\subsection{Gluon color structures beyond the diagonal core}

The diagonal color-density core contains the adjoint identity along the
active gluon line.  Wilson-line and gluonic-pole expansions introduce
additional adjoint indices.  The independent three-adjoint contractions are
\cite{Buffing2013,Boer2016}
\begin{equation}
 C_f^{abc}=if^{abc},
 \qquad
 C_d^{abc}=d^{abc}.
\label{eq:fd-color-tensors}
\end{equation}
They define distinct operator channels.  A process may probe
\begin{equation}
 W_{g,\mathrm{proc}}
 =C_f^{\mathrm{proc}}W_g^{(f)}
 +C_d^{\mathrm{proc}}W_g^{(d)},
\label{eq:process-fd-combination}
\end{equation}
but the coefficients belong to \(\Proc\), not to the universal parent
storage.

\begin{requirement}[Independent color members]
\label{req:independent-color-members}
The \(f\)- and \(d\)-type contributions retain separate state, link,
operator, uncertainty, and validation identities through the final hard
contraction.  A universal scalar ``gluon T-odd phase'' is forbidden.
\end{requirement}

\subsection{Gluon scalar-basis construction}

The complete leading-twist gluon GTMD basis for a spin-\(\tfrac12\) target is
constructed from the tensor product
\begin{equation}
 \left(0\oplus0\oplus2\right)_{\rm gluon\ transverse}
 \otimes
 \left(0\oplus1\right)_{\rm target\ helicity}
 \otimes
 \left\{\kT,\DT\right\}_{SO(2)},
\label{eq:gluon-basis-product}
\end{equation}
subject to Hermiticity, parity, and leading-twist power counting.  As in the
quark sector, a published notation is an adapter.  The production
implementation constructs candidate tensors, removes linear dependencies,
forms the Gram matrix, and generates dual projectors.

At \(\DT=0\), the scalar TMD basis for a spin-\(\tfrac12\) target contains
\begin{equation}
 \left\{
 f_1^g,
 h_1^{\perp g},
 g_{1L}^g,
 h_{1L}^{\perp g},
 f_{1T}^{\perp g},
 g_{1T}^g,
 h_{1T}^g,
 h_{1T}^{\perp g}
 \right\},
\label{eq:gluon-tmd-inventory}
\end{equation}
with the transverse ranks and link/color requirements recorded in
\cref{app:tmd-registry}.

\section{Helicity, transverse harmonics, and orbital selection}
\label{sec:helicity-oam}

\subsection{Joint parton--target helicity matrix}

For quarks or gluons, define
\begin{equation}
 \cM^a_{\lambda_a'\Lambda';\lambda_a\Lambda}
 (x,\kT,\DT)
\label{eq:joint-helicity-matrix}
\end{equation}
by contracting the operator with active-parton helicity polarization vectors
or spinors.  The nucleon target space has dimension two; the leading physical
parton-helicity space also has dimension two.  At fixed kinematics,
\(\cM^a\) is therefore a \(4\times4\) matrix before scalar decomposition.
Off forward it is not generally positive, but it remains the most direct
computational carrier of spin information.

\subsection{SO(2) weight and total angular momentum}

Introduce complex transverse components
\begin{equation}
 k_R=k_x+ik_y,
 \qquad
 k_L=k_x-ik_y,
 \qquad
 \Delta_R=\Delta_x+i\Delta_y,
 \qquad
 \Delta_L=\Delta_x-i\Delta_y.
\label{eq:complex-transverse}
\end{equation}
Under a transverse rotation by \(\phi\),
\begin{equation}
 k_R\mapsto e^{+i\phi}k_R,
 \qquad
 k_L\mapsto e^{-i\phi}k_L,
\label{eq:so2-weight}
\end{equation}
so a monomial carries an integer \(SO(2)\) weight.  For each amplitude block,
rotational equivariance imposes
\begin{equation}
 m_{\rm tensor}
 =\left(\Lambda'-\Lambda\right)
  +\left(\lambda_a'-\lambda_a\right)
  +\Delta L_z,
\label{eq:oam-selection-rule}
\end{equation}
up to the signed-helicity convention used to label the operator basis.  The
implementation derives the sign from the represented rotation generator; it
does not hard-code \cref{eq:oam-selection-rule} as a convention-independent
string.

\begin{proposition}[Interference zero test]
\label{prop:interference-zero-test}
Suppose a scalar GTMD projector has nonzero support only on interference
between two state blocks with orbital labels \(L_z\) and \(L_z'\).  Removing
either block makes that projected GTMD vanish.
\end{proposition}

\begin{proof}
The relevant matrix element is bilinear in the two blocks.  If either factor
is zero, every contraction selected by the projector is zero.
\end{proof}

\subsection{Physical interpretation of selected sectors}

The overlap formalism makes the following origins testable rather than
assigned:
\begin{itemize}
\item helicity-conserving rank-zero densities arise from diagonal spin/OAM
      blocks;
\item worm-gear and spin--orbit structures require helicity--orbital
      interference of the corresponding \(SO(2)\) weight;
\item pretzelosity and linearly polarized gluon structures require higher
      transverse rank and appropriate \(|\Delta L_z|\) content;
\item \(F_{1,4}\)- and \(G_{1,1}\)-type structures require nonzero
      \(\crossT{\kT}{\DT}\) and cannot be reconstructed from a forward scalar
      density;
\item naive-\(T\)-odd functions additionally require an absorptive link or
      rescattering phase, treated in Volume~III.
\end{itemize}

\begin{requirement}[Orbital-support manifest]
\label{req:orbital-support-manifest}
Every projected GTMD and TMD declares the pairs of target-helicity,
parton-helicity, and orbital blocks on which its projector has support.  This
manifest is used to generate zero tests and to diagnose missing basis sectors.
\end{requirement}

\section{Discrete symmetries and link reversal}
\label{sec:symmetries}

\subsection{Hermiticity}

For the bilocal quark operator,
\begin{equation}
 \left[\bar\psi(-z/2)\Gamma U_\gamma\psi(z/2)\right]^\dagger
 =\bar\psi(z/2)\,
 \gamma^0\Gamma^\dagger\gamma^0\,
 U_{\gamma^{-1}}\psi(-z/2).
\label{eq:bilocal-dagger}
\end{equation}
Consequently,
\begin{align}
 \left[W_{\Lambda'\Lambda}^{[\Gamma]}
 (x,\kT,\DT;\gamma)\right]^*
 &=W_{\Lambda\Lambda'}^{[\gamma^0\Gamma^\dagger\gamma^0]}
 (x,\kT,-\DT;\gamma^{-1}),
\label{eq:quark-hermiticity}\\
 \left[W_{\Lambda'\Lambda}^{g,ij}
 (x,\kT,\DT;\gamma_1,\gamma_2)\right]^*
 &=W_{\Lambda\Lambda'}^{g,ji}
 (x,\kT,-\DT;\gamma_2^{-1},\gamma_1^{-1}).
\label{eq:gluon-hermiticity}
\end{align}
The exact endpoint and link ordering in the second relation is part of the
operator adapter.  For the leading Hermitian quark projectors, the Dirac
label returns to itself.

\begin{proposition}[Overlap Hermiticity]
\label{prop:overlap-hermiticity}
If the state phase convention is common, the active kernels satisfy the
adjoint relations above, and the recoil map is
\cref{def:recoil-map}, then the canonical overlaps obey
\cref{eq:quark-hermiticity,eq:gluon-hermiticity}.
\end{proposition}

\begin{proof}
Complex conjugation interchanges the incoming and outgoing wave functions,
conjugates the active kernel, and reverses the operator path.  By
\cref{prop:recoil-jacobian}, \(\DT\to-\DT\) interchanges the two recoil
coordinate sets with no Jacobian.  Relabeling dummy spectator variables gives
the stated relations.
\end{proof}

\subsection{Parity}

Light-front parity is implemented as a represented transformation on the
momentum fibers, target helicities, parton helicities, transverse tensors,
and operator basis.  We denote it by
\begin{equation}
 \cU_P:\Hil_{\mathrm{phys}}(p;r)\to\Hil_{\mathrm{phys}}(p_P;r),
\label{eq:lf-parity-map}
\end{equation}
with its exact action fixed by the selected front-form convention.  The
operator basis must satisfy
\begin{equation}
 \cU_P\cO_{\Gamma,\gamma}(p',p)\cU_P^{-1}
 =\pi_\Gamma\cO_{\Gamma_P,\gamma_P}(p'_P,p_P).
\label{eq:parity-operator}
\end{equation}
Parity reduces the number of independent helicity amplitudes but does not
justify identifying different flavors, links, or transverse ranks.

\subsection{Time reversal and process links}

Time reversal is antiunitary and maps future- and past-pointing staples into
one another.  For a scalar projection \(F\) with declared naive-time parity
\(\tau_F=\pm1\),
\begin{equation}
 F^{[\gamma_+]}(x,\kT,\DT)
 =\tau_F\,
 \mathfrak T_F
 \left[F^{[\gamma_-]}(x,-\kT,-\DT)\right]^*,
\label{eq:time-reversal-link}
\end{equation}
where \(\mathfrak T_F\) is the helicity/tensor phase adapter.  In the forward
TMD limit and a common convention, this yields the familiar relations
\begin{equation}
 F_{\rm even}^{[+]}=F_{\rm even}^{[-]},
 \qquad
 F_{\rm odd}^{[+]}=-F_{\rm odd}^{[-]}.
\label{eq:tmd-link-reversal}
\end{equation}

\begin{requirement}[Operator-level reversal]
\label{req:operator-level-reversal}
Link reversal is applied to the decorated operator before projection.  A
sign may not be attached downstream to a named function without verifying
its operator, color, and tensor conventions.
\end{requirement}

\subsection{Real-overlap limit}

If all retained amplitudes, active kernels, and zeroth-order links are real in
a common phase convention, every projection that is odd under complex
conjugation vanishes.  This is a controlled zero-rescattering limit, not a
claim that the physical Sivers, Boer--Mulders, or gluon T-odd functions are
zero.

\begin{proposition}[Zero-phase test]
\label{prop:zero-phase-test}
Let \(\cP_A\) select the imaginary part of an interference block.  If the
participating amplitudes and kernel are real, then \(F_A=0\).
\end{proposition}

\begin{proof}
The selected bilinear has the form
\(\operatorname{Im}(\psi_a^*K\psi_b)\).  Under the hypothesis it is real, so
its imaginary part vanishes.
\end{proof}

\section{Typed reductions from the common parent}
\label{sec:reductions}

\subsection{The reduction diagram}

The central reduction architecture is
\begin{equation}
\begin{tikzcd}[column sep=large,row sep=large]
 W_{a/N,\LF}^{(r)}(x,\kT,\DT;\gamma)
 \arrow[r,"\cZ_{\LF\to\rm QCD}"]
 \arrow[d,"\mathsf T_{\rm reg}"']
 \arrow[dr,phantom,"\scriptstyle\circlearrowright"]
 & W_{a/N,\rm QCD}(x,\kT,\DT;\mu,\zeta,\mathfrak s)
 \arrow[d,"\mathsf T_{\rm QCD}"]
 \arrow[r,"\mathsf W_\Delta"]
 & \rho_{a/N}(x,\kT,\bD;\mu,\zeta,\mathfrak s)
 \\
 \Phi_{a/N,\LF}^{(r)}(x,\kT)
 \arrow[r,"\cZ_{\LF\to\TMD}"]
 & \Phi_{a/N,\TMD}(x,\kT;\mu,\zeta,\mathfrak s)
 \arrow[r,"\mathsf P_A"]
 & F_{a/N,A}(x,\kT;\mu,\zeta)
\end{tikzcd}
\label{eq:tmd-commuting-diagram}
\end{equation}
A second matched route gives collinear and local operators:
\begin{equation}
\begin{tikzcd}[column sep=large,row sep=large]
 W_{a/N,\LF}^{(r)}
 \arrow[r,"\cZ_{\LF\to\rm QCD}"]
 \arrow[d,"\mathsf G_{\rm reg}"']
 & W_{a/N,\rm QCD}
 \arrow[d,"\mathsf M_{\rm link/OPE}"]
 \\
 H_{a/N,\LF}^{(r)}(x,0,t)
 \arrow[r,"\cZ_{\LF\to\GPD}"]
 & H_{a/N,\GPD}(x,0,t;\mu)
 \arrow[r,"\int\dd x\,x^n"]
 & \langle p'|\cO_{n}^{\rm local}(\mu)|p\rangle
\end{tikzcd}
\label{eq:gpd-commuting-diagram}
\end{equation}
The curved-arrow symbol denotes a required commutativity test up to the
combined matching, truncation, and numerical error.

\subsection{TMD forward limit}

The regulated forward map is
\begin{equation}
 \mathsf T_{\rm reg}[W]
 =W(x,\kT,\DT=0;\gamma),
\label{eq:tmd-forward-map}
\end{equation}
with target and active-parton helicity indices retained.  The physical TMD is
not merely \cref{eq:tmd-forward-map}; it is the corresponding UV-renormalized,
rapidity-renormalized, soft-subtracted operator in the chosen scheme.

\begin{requirement}[Path-preserving forward limit]
\label{req:path-preserving-forward}
Setting \(\DT=0\) does not erase the Wilson path, ordered gluon-link pair,
rapidity regulator, soft prescription, or process-link class.
\end{requirement}

\subsection{Partonic Wigner transform}

The transverse Wigner coordinate is \(\bD\), conjugate to \(\DT\):
\begin{equation}
 \rho_{a/N}^{[J]}(x,\kT,\bD;\gamma)
 =\int\frac{\dd^2\DT}{(2\pi)^2}
 e^{-i\bD\cdot\DT}
 W_{a/N}^{[J]}(x,\kT,\DT;\gamma).
\label{eq:wigner-transform}
\end{equation}
The TMD impact separation \(\bM\), conjugate to \(\kT\), is a different
coordinate:
\begin{equation}
 \widetilde F(x,\bM)
 =\int\dd^2\kT\,e^{+i\bM\cdot\kT}F(x,\kT).
\label{eq:tmd-fourier-transform}
\end{equation}
Equal units do not permit these objects to be aliased.

\subsection{GPD route and link matching}

A formal regulated \(\kT\) integral sets the transverse operator separation
to zero:
\begin{equation}
 \mathsf G_{\rm reg}[W]
 =\int\dd^2\kT\,W(x,\kT,\DT;\gamma).
\label{eq:regulated-gpd-integral}
\end{equation}
For a straight collinear link and a compatible regulator, this becomes a
regulated GPD matrix element.  For a staple operator with rapidity and soft
subtraction, the renormalized GPD route contains an explicit link-shortening,
soft, and ultraviolet matching map:
\begin{equation}
 H_{\rm GPD}(x,0,t;\mu)
 =\mathsf M_{\rm link/OPE}(\mu,\zeta,\mathfrak s)
 \left[\int\dd^2\kT\,
 W_{\rm GTMD}(x,\kT,\DT;\mu,\zeta,\mathfrak s)\right].
\label{eq:matched-gpd-route}
\end{equation}
The same warning applies to a literal integral of a soft-subtracted TMD.

\begin{requirement}[No naive staple integral]
\label{req:no-naive-staple-integral}
A soft-subtracted staple GTMD/TMD may not be identified with a GPD, PDF, or
local current solely by numerical integration.  The operator path and
renormalization must first be matched to the relevant collinear operator.
\end{requirement}

\subsection{PDF and local-current routes}

At zero transfer, the matched collinear distribution is
\begin{equation}
 f_{a/N}^{[J]}(x;\mu)
 =\mathsf P_{\rm PDF}
 \mathsf M_{\rm link/OPE}
 \mathsf G[W_{a/N}^{[J]}](x,0;\mu).
\label{eq:pdf-route}
\end{equation}
Local moments are
\begin{equation}
 \int_{-1}^{1}\dd x\,x^n H^{[J]}(x,\xi,t;\mu)
 =\sum_k C_{nk}^{[J]}(\xi,t;\mu)
 \langle p'|\cO_{nk}^{[J]}(\mu)|p\rangle,
\label{eq:local-moment-route}
\end{equation}
with the polynomiality structure becoming fully testable only when the
\(\xi\ne0\) extension is implemented.

\begin{theorem}[Matched common-parent consistency]
\label{thm:matched-common-parent}
Let all four paths in \cref{eq:tmd-commuting-diagram,eq:gpd-commuting-diagram}
be defined in compatible regulator and operator schemes.  If their matching
coefficients are derived to accuracy \(\mathcal A\) and the state is converged
to truncation accuracy \(\delta_r\), then two routes to the same renormalized
matrix element differ by
\begin{equation}
 \Delta\cO
 =\cO(\delta_r)+\cO(\delta_{\mathcal A})
 +\cO(\delta_{\rm num}).
\label{eq:commuting-error}
\end{equation}
A discrepancy larger than the combined error falsifies at least one state,
operator, recoil, matching, or numerical component.
\end{theorem}

\begin{proof}
Each path is a composition of linear regulated matrix-element maps and
matching morphisms representing the same renormalized target operator up to
their stated truncations.  Subtracting the two compositions cancels the
common exact contribution, leaving the state, matching, and numerical
remainders.  The conclusion follows from the triangle inequality for the
chosen observable norm.
\end{proof}

\section{Projection to named nucleon TMDs}
\label{sec:tmd-projection}

\subsection{Projector definition}

At \(\DT=0\), named TMDs are linear coordinates of the joint helicity matrix:
\begin{equation}
 F_{a/N,A}(x,\kT;\mathfrak o)
 =\Tr_{\lambda_a,\Lambda}
 \left[
 \cP_A(\widehat\kT,S)\,
 \cM^a(x,\kT,0;\mathfrak o)
 \right].
\label{eq:named-tmd-projector}
\end{equation}
The projector identity contains
\begin{equation}
 \begin{aligned}
 \mathfrak p_A=\bigl(&A,a,J,\text{target polarization},
 \text{parton polarization},m,n_k,n_b,\\
 &M_{\rm ref},J_{|m|},\phi_F,
 \text{link parity},\text{collinear status}\bigr).
 \end{aligned}
\label{eq:tmd-projector-id}
\end{equation}

\begin{requirement}[Named functions are views]
\label{req:named-functions-views}
A named TMD is a view produced by \cref{eq:named-tmd-projector}; it is not an
independent mutable field of the microscopic state.  Caching is allowed only
with the complete parent, projector, scheme, scale, and member identity.
\end{requirement}

\subsection{Rank-aware Fourier--Bessel transform}

For an irreducible harmonic of weight \(m\), define
\begin{align}
 \widetilde F^{(m)}(x,b)
 &=2\pi i^m
 \int_0^\infty k\,\dd k\,
 J_{|m|}(bk)\,N_m(k;M_{\rm ref})F^{(m)}(x,k),
\label{eq:rank-forward-transform}\\
 F^{(m)}(x,k)
 &=\frac{(-i)^m}{2\pi}
 \int_0^\infty b\,\dd b\,
 J_{|m|}(bk)\,\widetilde N_m(b;M_{\rm ref})
 \widetilde F^{(m)}(x,b).
\label{eq:rank-inverse-transform}
\end{align}
The convention factors \(N_m\), \(\widetilde N_m\), and the phase \(i^m\)
are stored rather than inferred.  A scalar \(J_0\) routine is invalid for a
rank-one or rank-two coefficient even when its input is a scalar array.

\begin{proposition}[Angular integral]
\label{prop:angular-integral}
An uncontracted irreducible transverse harmonic of nonzero weight has zero
unweighted azimuthal integral.
\end{proposition}

\begin{proof}
The angular dependence is proportional to \(e^{im\phi}\) with
\(m\ne0\), and
\(\int_0^{2\pi}\dd\phi\,e^{im\phi}=0\).
\end{proof}

\subsection{Collinear-status classes}

Each TMD is assigned one of:
\begin{enumerate}
\item \texttt{DIRECT}: a nonzero rank-zero collinear operator limit;
\item \texttt{MOMENT\_ONLY}: zero unweighted integral but a defined weighted
      transverse moment;
\item \texttt{MATCHING\_GENERATED}: no independent leading collinear density,
      but a perturbative small-\(\bM\) coefficient;
\item \texttt{NO\_COLLINEAR\_LIMIT}: the relevant integrated operator is zero
      or undefined in the declared channel;
\item \texttt{UNRESOLVED}: matching or moment status remains to be derived.
\end{enumerate}
This status belongs to the projector/operator convention, not to the
numerical appearance of a curve.

\section{GPDs, form factors, and local sum rules}
\label{sec:gpd-moments}

\subsection{Vector and axial GPDs}

After the matched \(\kT\) integration, the standard spin-\(\tfrac12\) vector
and axial decompositions are
\begin{align}
 \frac{1}{2P^+}\bar u(p',\Lambda')
 \left[\gamma^+H^q(x,\xi,t)
 +\frac{i\sigma^{+\alpha}\Delta_\alpha}{2M_N}E^q(x,\xi,t)
 \right]u(p,\Lambda),
\label{eq:vector-gpd-decomp}\\
 \frac{1}{2P^+}\bar u(p',\Lambda')
 \left[\gamma^+\gamma_5\widetilde H^q(x,\xi,t)
 +\frac{\Delta^+\gamma_5}{2M_N}\widetilde E^q(x,\xi,t)
 \right]u(p,\Lambda).
\label{eq:axial-gpd-decomp}
\end{align}
At \(\xi=0\), the \(\widetilde E\) term in this particular plus-component
projection is kinematically suppressed, but the function remains part of the
full nonzero-skewness operator family.

The chiral-odd GPD sector contains
\begin{equation}
 H_T^q,
 \quad E_T^q,
 \quad\widetilde H_T^q,
 \quad\widetilde E_T^q,
\label{eq:chiral-odd-gpds}
\end{equation}
and is obtained from the matched integral of the same
\(i\sigma^{j+}\gamma_5\) parent.

\subsection{Electromagnetic and axial moments}

With standard local-current normalization,
\begin{align}
 \int_{-1}^{1}\dd x\,H^q(x,\xi,t;\mu)&=F_1^q(t),
\label{eq:h-f1-sum}\\
 \int_{-1}^{1}\dd x\,E^q(x,\xi,t;\mu)&=F_2^q(t),
\label{eq:e-f2-sum}\\
 \int_{-1}^{1}\dd x\,\widetilde H^q(x,\xi,t;\mu)&=G_A^q(t),
\label{eq:htilde-ga-sum}
\end{align}
subject to the current-renormalization convention.  The right-hand sides are
independent of \(\xi\); failure of this property in a nonzero-skewness
extension is a polynomiality or current-consistency defect.

\subsection{Energy--momentum tensor and spin}

The Mellin power depends on the species normalization.  With the standard
quark convention, the second vector moments have the schematic polynomiality
form
\begin{align}
 \int_{-1}^{1}\dd x\,xH^q(x,\xi,t;\mu)
 &=A_q(t;\mu)+(2\xi)^2 C_q(t;\mu)+\cdots,
\label{eq:emt-h-moment}\\
 \int_{-1}^{1}\dd x\,xE^q(x,\xi,t;\mu)
 &=B_q(t;\mu)-(2\xi)^2 C_q(t;\mu)+\cdots.
\label{eq:emt-e-moment}
\end{align}
For the common gluon convention in which
\(H^g(x,0,0)=xg(x)\), the corresponding energy--momentum moments are
\begin{align}
 \int_{0}^{1}\dd x\,H^g(x,\xi,t;\mu)
 &=A_g(t;\mu)+(2\xi)^2 C_g(t;\mu)+\cdots,
\label{eq:emt-gluon-h-moment}\\
 \int_{0}^{1}\dd x\,E^g(x,\xi,t;\mu)
 &=B_g(t;\mu)-(2\xi)^2 C_g(t;\mu)+\cdots.
\label{eq:emt-gluon-e-moment}
\end{align}
Factors multiplying the \(C_a\) term vary with polynomiality convention and
are therefore adapter data; the invariant requirement is that one declared
Mellin convention be used consistently in the GTMD, GPD, and local-operator
routes.  At \(t=0\), the kinetic angular momentum relation is \cite{Ji1997}
\begin{equation}
 J_a(\mu)=\frac12\left[A_a(0;\mu)+B_a(0;\mu)\right],
 \qquad a=q,g.
\label{eq:ji-relation}
\end{equation}
The sum over quarks and gluons must satisfy the matched momentum and angular-
momentum ledgers.  Canonical OAM from a staple-link Wigner operator need not
equal the kinetic OAM in \cref{eq:ji-relation}; their difference is an
operator and gauge-link question, not a numerical discrepancy to be tuned
away.

\section{Wigner distributions, OAM, and spin--orbit moments}
\label{sec:wigner-oam}

\subsection{Phase-space quasidistribution}

For a fixed operator identity, the partonic Wigner transform is
\cref{eq:wigner-transform}.  Its natural transverse phase space is
\begin{equation}
 \mathcal P_T=T^*\mathbb R_T^2
 =\{(\bD,\kT)\},
 \qquad
 \omega=\dd b_{\Delta x}\wedge\dd k_x
       +\dd b_{\Delta y}\wedge\dd k_y.
\label{eq:transverse-phase-space}
\end{equation}
The rotation moment map is
\begin{equation}
 \mu_{L_z}(\bD,\kT)
 =b_{\Delta x}k_y-b_{\Delta y}k_x.
\label{eq:oam-moment-map}
\end{equation}
The Wigner distribution is a quasidistribution and need not be positive
pointwise \cite{BelitskyJiYuan2004,LorcePasquini2011}.

\subsection{Canonical phase-space OAM}

For the unpolarized active-quark projection in a longitudinally polarized
nucleon, define the link-dependent phase-space OAM
\begin{equation}
 L_z^q[\gamma]
 =\int\dd x\,\dd^2\kT\,\dd^2\bD\,
 (\bD\times\kT)_z\,
 \rho_{LU}^q(x,\kT,\bD;\gamma).
\label{eq:wigner-oam}
\end{equation}
The operator path determines whether the quantity is connected to a straight-
link kinetic-type definition or a staple-link canonical/Jaffe--Manohar-type
definition, with the precise relation depending on boundary and gauge
conventions.

In a commonly used \(F_{1,4}\) adapter, the same moment is represented by
\cite{LorcePasquini2011,Kanazawa2014}
\begin{equation}
 L_z^q[\gamma]
 =-\int\dd x\,\dd^2\kT\,
 \frac{\kT^2}{M_N^2}
 F_{1,4}^q(x,0,\kT^2,0,0;\gamma),
\label{eq:f14-oam-adapter}
\end{equation}
where the sign, mass factor, and link are part of the adapter.  The production
code verifies \cref{eq:f14-oam-adapter} from the Fourier and helicity
conventions rather than treating it as an untyped identity.

\subsection{Spin--orbit correlation}

A longitudinal quark-spin--orbital correlation can analogously be defined
for an unpolarized nucleon and a longitudinally polarized quark by the axial
Wigner moment,
\begin{equation}
 C_z^q[\gamma]
 =\int\dd x\,\dd^2\kT\,\dd^2\bD\,
 (\bD\times\kT)_z\,
 \rho_{UL}^q(x,\kT,\bD;\gamma),
\label{eq:spin-orbit-wigner}
\end{equation}
and in a standard \(G_{1,1}\) adapter has the form
\begin{equation}
 C_z^q[\gamma]
 =\int\dd x\,\dd^2\kT\,
 \frac{\kT^2}{M_N^2}
 G_{1,1}^q(x,0,\kT^2,0,0;\gamma),
\label{eq:g11-spin-orbit}
\end{equation}
again subject to the declared convention.

\begin{requirement}[OAM operator identity]
\label{req:oam-operator-identity}
Every OAM or spin--orbit number stores the Wilson path, gauge/boundary
prescription, renormalization scheme, scale, and distinction between
canonical and kinetic operators.  A plot label ``quark OAM'' is insufficient.
\end{requirement}

\subsection{Phase-space closure tests}

The Wigner transform must satisfy
\begin{align}
 \int\dd^2\bD\,
 \rho(x,\kT,\bD)&=W(x,\kT,\DT=0),
\label{eq:wigner-tmd-marginal}\\
 \int\dd^2\kT\,
 \rho(x,\kT,\bD)&=q(x,\bD)
 \quad\text{after the required collinear link matching},
\label{eq:wigner-impact-marginal}\\
 \int\dd^2\bD\,\dd^2\kT\,
 \rho(x,\kT,\bD)&=f(x)
 \quad\text{in a compatible scheme}.
\label{eq:wigner-pdf-marginal}
\end{align}
The first relation is a direct Fourier identity.  The latter two include the
same operator-matching qualifications as \cref{sec:reductions}.

\section{Positivity and off-forward bounds}
\label{sec:positivity}

\subsection{Forward positivity domain}

The Gram representation of \cref{prop:forward-gram} implies
\begin{equation}
 \Phi_a(x,\kT)=\Phi_a^\dagger(x,\kT),
 \qquad
 v^\dagger\Phi_a(x,\kT)v\ge0
\label{eq:forward-psd}
\end{equation}
for the applicable forward amplitude-level density correlator.  Named TMD
bounds follow from principal minors or semidefinite optimization in the full
joint parton--nucleon helicity matrix.

\begin{requirement}[No positivity clipping]
\label{req:no-positivity-clipping}
A negative eigenvalue may not be repaired by clipping the assembled matrix.
It diagnoses an inconsistent state, active kernel, projector, interpolation,
or numerical tolerance.  Any modified amplitude must be propagated through
all common-parent reductions.
\end{requirement}

\subsection{Off-forward Cauchy--Schwarz bound}

Although an off-forward GTMD matrix is not positive semidefinite, the overlap
obeys an operator-norm Cauchy--Schwarz inequality.  For fixed external
indices, write
\begin{equation}
 W_{\beta\alpha}(\Delta)
 =\langle A_{\beta}^{\rm out},
 \cK_{\beta\alpha}A_{\alpha}^{\rm in}\rangle,
\label{eq:offforward-amplitude-form}
\end{equation}
where the Hilbert-space index includes sector, spectator, momentum, color,
and helicity labels.  Then
\begin{equation}
 |W_{\beta\alpha}(\Delta)|^2
 \le
 \norm{\cK_{\beta\alpha}}_{\rm op}^{2}
 \norm{A_{\beta}^{\rm out}}^2
 \norm{A_{\alpha}^{\rm in}}^2.
\label{eq:offforward-cs-bound}
\end{equation}
For an elementary normalized helicity selector or unitary spin transition,
\(\norm{\cK}_{\rm op}=1\), recovering the familiar product of recoil-shifted
diagonal densities.  The right-hand side is evaluated at the appropriate
incoming and outgoing recoil kinematics, not necessarily at the same
\((x,\kT)\) point as a forward TMD.

\begin{proposition}[Off-forward overlap bound]
\label{prop:offforward-bound}
The canonical overlap satisfies \cref{eq:offforward-cs-bound} sector by
sector and after the orthogonal direct sum over retained sectors.
\end{proposition}

\begin{proof}
Apply Cauchy--Schwarz to
\(\langle A^{\rm out},\cK A^{\rm in}\rangle\), followed by the operator-norm
bound \(\norm{\cK A^{\rm in}}\le
\norm{\cK}_{\rm op}\norm{A^{\rm in}}\).  Orthogonal direct sums preserve the
same inequality with the block-operator norm.
\end{proof}

\subsection{Where exact positivity is not imposed}

Exact pointwise PSD is not required for:
\begin{itemize}
\item off-forward GTMD matrices;
\item Wigner quasidistributions;
\item scalar GTMD coefficients separated from their complete helicity matrix;
\item finite-order, UV-renormalized, rapidity-subtracted scheme objects unless
      a positivity theorem has been established for that definition;
\item distributions after truncating perturbative matching in a way that is
      not positivity preserving.
\end{itemize}
The model stores which layer carries a Gram proof and which carries only
matched or perturbative consistency.

\section{Regulator, basis, and Fock-sector convergence}
\label{sec:convergence}

\subsection{Common truncation tower}

Let \(r\preceq r'\) denote a refinement in basis resolution, ultraviolet or
infrared regulator, Fock content, quadrature, or tensor-network rank.  A
common-parent observable at level \(r\) is
\begin{equation}
 \cO_A^{(r)}
 =\mathsf R_A\cZ_r
 \langle N;r|\cO_r|N;r\rangle.
\label{eq:observable-at-r}
\end{equation}
Convergence is tested after matching:
\begin{equation}
 \delta_A(r',r)
 =\left\|
 \mathsf R_{r'\to r}\!\left[\cO_A^{(r')}\right]
 -\cO_A^{(r)}
 \right\|.
\label{eq:gtmd-convergence-residual}
\end{equation}
No single smooth curve at one cutoff establishes convergence.

\subsection{Required axes}

The GTMD convergence manifest reports separately:
\begin{enumerate}
\item longitudinal resolution and endpoint treatment;
\item transverse basis/grid and ultraviolet support;
\item infrared regulator and large-distance basis;
\item retained Fock sectors, especially \(qqqg\) and \(qqqq\bar q\);
\item orbital and helicity blocks;
\item Hamiltonian and operator counterterm order;
\item Wilson/eikonal order when enabled;
\item recoil and interpolation discretization;
\item \(\DT\), \(\kT\), and Fourier quadrature;
\item tensor-network or low-rank contraction accuracy.
\end{enumerate}

\begin{requirement}[No cross-category hiding]
\label{req:no-cross-category-hiding}
Agreement under grid refinement cannot substitute for Fock-sector
convergence; positivity cannot substitute for regulator convergence; and a
small posterior band cannot substitute for matching uncertainty.  Each axis
is reported independently before any combined covariance is constructed.
\end{requirement}

\subsection{Common-state covariance}

For microscopic parameters \(\theta\) and controlled truncation choices
\(\nu\),
\begin{equation}
 \operatorname{Cov}(\cO_A,\cO_B)
 =\left\langle
 \left(\cO_A-\langle\cO_A\rangle\right)
 \left(\cO_B-\langle\cO_B\rangle\right)
 \right\rangle_{\theta,\nu}.
\label{eq:common-parent-covariance}
\end{equation}
The same ensemble member is propagated through quark, antiquark, gluon,
TMD, GPD, current, and OAM reductions.  Independent resampling of each named
function would destroy the predictive correlations.

\section{Computational object model}
\label{sec:software}

\subsection{Core interfaces}

\begin{longtable}{L{0.25\textwidth}L{0.65\textwidth}}
\toprule
Object & Formal responsibility\\
\midrule
\endhead
\texttt{MomentumFiber} & Total hadron momentum, LF convention, on-shell mass, regulator, state-normalization identity, and compatible physical Hilbert space.\\
\texttt{RecoilMap} & Versioned incoming/outgoing intrinsic coordinates, active index, transfer type, Jacobian, inverse, and closure tests.\\
\texttt{ActivePartonSelector} & Species, flavor, active slot, permutation orbit, occupation multiplicity, and spectator identity.\\
\texttt{OverlapKernel} & Active spin/color operator, spectator deltas, field normalization, Wilson order, source and target sectors, and adjoint.\\
\texttt{QuarkGTMDOperator} & Decorated quark bilocal, Dirac projection, link, regulator/scheme data, and positive/negative-\(x\) representation.\\
\texttt{AntiquarkGTMDAdapter} & Positive-\(x\) active-antiquark evaluator and explicit charge-conjugation conversion.\\
\texttt{GluonGTMDOperator} & Field-strength indices, ordered adjoint link pair, color-contraction status, polarization projectors, and normalization anchor.\\
\texttt{GTMDStateMatrix} & Joint target--parton helicity matrix with kinematics, microscopic member, operator identity, and convergence metadata.\\
\texttt{GTMDProjector} & Basis tensor, dual projector, scalar name, rank/mass convention, symmetry behavior, and reference-convention adapter.\\
\texttt{ReductionMap} & Typed TMD, Wigner, GPD, PDF, moment, and local-current routes with matching prerequisites.\\
\texttt{CommonParent\allowbreak Closure\allowbreak Report} & Route-by-route residuals, tolerance budget, marginal/current/OAM checks, and provenance path.\\
\texttt{GTMD\allowbreak Convergence\allowbreak Manifest} & Regulator, basis, sector, OAM, Wilson, quadrature, and matching sequences with observable residuals.\\
\bottomrule
\end{longtable}

\subsection{Immutability and cache keys}

A cached scalar function is valid only under the key
\begin{equation}
 \mathfrak K=
 (\text{microscopic member},\mathfrak o,\mathfrak R_{j,\Delta},
 \mathfrak p_A,\text{reduction path},\mu,\zeta,
 \text{grid/quadrature manifest}).
\label{eq:gtmd-cache-key}
\end{equation}
Changing any field invalidates the cache.  Mutable global conventions are
forbidden.

\subsection{Composition predicate}

For typed maps \(f:A\to B\) and \(g:C\to D\), composition is permitted only
if
\begin{equation}
 \operatorname{Compatible}(B,C)
 =\texttt{true},
\label{eq:composition-predicate}
\end{equation}
where compatibility includes momentum fiber, coordinate role, operator
scheme, link class, color representation, rank convention, state member,
normalization, and truncation.  Otherwise the code must return a structured
mismatch report listing every nonmatching field.

\begin{requirement}[No permissive fallback]
\label{req:no-permissive-fallback}
The implementation may not coerce an unresolved field to a convenient
default, strip metadata to make two objects compare equal, or silently
recompute a scalar with a different projector.
\end{requirement}

\subsection{Stable formal requirement identifiers}

\begin{longtable}{L{0.16\textwidth}L{0.34\textwidth}L{0.40\textwidth}}
\toprule
ID & Formal source & Implementation obligation\\
\midrule
\endhead
V2.KIN & \cref{sec:kinematics,sec:recoil} & Implement momentum fibers and the unique symmetric \(\xi=0\) recoil authority.\\
V2.OPID & \cref{def:decorated-operator} & Carry complete quark/gluon operator and ordered-link identity.\\
V2.OVERLAP & \cref{eq:master-diagonal-overlap} & Evaluate a common state with explicit active and spectator kernels.\\
V2.OFFDIAG & \cref{req:named-off-diagonal} & Reject number-changing blocks without a named source.\\
V2.QUARK & \cref{sec:quark} & Implement good-field quark kernels and sixteen-function projector basis.\\
V2.ANTIQ & \cref{sec:antiquark} & Use positive-\(x\) antiquark sectors and explicit charge conjugation.\\
V2.GLUON & \cref{sec:gluon} & Implement field-strength overlap, ordered links, polarization tensors, and color status.\\
V2.REDUCE & \cref{sec:reductions} & Implement matched TMD/Wigner/GPD/PDF/current routes.\\
V2.RANK & \cref{sec:tmd-projection} & Use rank-aware transforms and collinear-status metadata.\\
V2.OAM & \cref{sec:wigner-oam} & Preserve \(\bD\) versus \(\bM\), link identity, and OAM convention adapters.\\
V2.POS & \cref{sec:positivity} & Prove forward Gram PSD; do not clip or overextend positivity.\\
V2.CONV & \cref{sec:convergence} & Report observable convergence on every independent truncation axis.\\
\bottomrule
\end{longtable}

\section{Property-based tests and validation ladder}
\label{sec:tests}

\subsection{Exact algebraic and kinematic tests}

The following tests are exact within floating-point or symbolic tolerance:
\begin{enumerate}
\item \textbf{LF invariant:}
      verify \(p^2=2p^+p^- -\bm p_T^2\) for every momentum fiber.
\item \textbf{Intrinsic closure:}
      verify \(\sum_i x_i=1\) and \(\sum_i\kiT=0\) before and after recoil.
\item \textbf{Physical recoil:}
      active physical momentum changes by \(\DT\); spectators do not change.
\item \textbf{Recoil inverse:}
      applying the reversed transfer interchanges incoming and outgoing sets.
\item \textbf{Unit Jacobian:}
      symbolic or automatic differentiation gives determinant one.
\item \textbf{Forward identity:}
      incoming and outgoing coordinates coincide at \(\DT=0\).
\item \textbf{Spectator identity:}
      every unresolved spectator label is matched exactly.
\item \textbf{Permutation equivalence:}
      explicit active-slot summation agrees with occupation-multiplicity form.
\item \textbf{Operator adjoint:}
      quark and gluon kernels satisfy their declared dagger relations.
\item \textbf{Hermiticity:}
      the full overlap obeys
      \cref{eq:quark-hermiticity,eq:gluon-hermiticity}.
\item \textbf{Projector duality:}
      \(\langle\cP_A,B_B\rangle=\delta_{AB}\) at non-degenerate kinematics.
\item \textbf{Rank covariance:}
      each projected scalar transforms with its declared \(SO(2)\) weight.
\item \textbf{Link identity:}
      path reversal, ordered gluon-link exchange, and color status are tested
      independently.
\item \textbf{Map mismatch injection:}
      coordinate, rank, mass, fiber, link, scheme, and map-class mismatches are
      rejected with structured diagnostics.
\end{enumerate}

\subsection{State and operator tests}

\begin{enumerate}
\item the same microscopic member ID appears in every species parent;
\item active flavor counting reproduces exact valence charges;
\item disabling sea-containing and induced sectors makes positive-\(x\)
      antiquarks vanish;
\item disabling gluon-containing and induced sectors makes the gluon parent
      vanish;
\item target and parton helicity blocks reconstruct the original correlator;
\item removing an OAM block removes every projector requiring its
      interference;
\item real-amplitude/zero-rescattering mode removes imaginary T-odd
      projections;
\item the vector current at zero transfer reproduces the state charge;
\item the matched quark-plus-gluon momentum moment closes within its
      truncation budget;
\item forward Gram matrices are PSD without eigenvalue clipping.
\end{enumerate}

\subsection{Reduction and closure tests}

\begin{enumerate}
\item direct \(\DT=0\) evaluation agrees with the TMD reduction map;
\item Fourier integration over \(\bD\) returns the forward GTMD;
\item rank-aware \(\kT\leftrightarrow\bM\) transforms round trip;
\item regulated \(\kT\) integration agrees with the regulated collinear
      overlap before scheme matching;
\item matched GPD moments reproduce currents computed directly with the
      Volume~I current operator;
\item independent routes to PDF, form-factor, OAM, and spin--orbit moments
      satisfy \cref{eq:commuting-error};
\item future/past link reversal holds for every projector with declared
      naive-time parity;
\item \(f\)- and \(d\)-type gluon parents remain separate until an explicit
      process contraction.
\end{enumerate}

\subsection{Validation ladder}

A nucleon parent advances through the following gates:
\begin{enumerate}
\item \textbf{Typed-object gate:} every coordinate, rank, operator, link,
      fiber, and scheme field is complete for the requested operation.
\item \textbf{Algebra gate:} recoil, adjoint, permutation, color, and projector
      identities pass.
\item \textbf{State gate:} normalization, charges, spectrum, and state
      convergence satisfy Volume~I.
\item \textbf{Parent gate:} quark, antiquark, and gluon overlaps are generated
      from the same member with exact support and Hermiticity.
\item \textbf{Marginal gate:} TMD, Wigner, GPD, PDF, and current routes close
      within the declared error budget.
\item \textbf{Physical gate:} withheld flavor, form-factor, helicity, OAM, or
      lattice moments are predicted without named-function retuning.
\item \textbf{Wilson gate:} after Volume~III, absorptive phases, link reversal,
      and \(f/d\) color sectors pass.
\item \textbf{Nuclear gate:} after Volume~IV, the same nucleon parent enters
      the spin-1 state without unrelated renormalization.
\end{enumerate}

\section{Analytic benchmark hierarchy}
\label{sec:benchmarks}

\subsection{Benchmark A: one-body point state}

For a one-constituent state, \(x_1=1\) and \(\bm{k}_{1T}=0\).  The recoil shifts
vanish:
\begin{equation}
 \bm\kappa_{1T}^{\rm in}
 =\bm\kappa_{1T}^{\rm out}=0.
\label{eq:one-body-recoil}
\end{equation}
The vector overlap reduces to the active current matrix element.  This
benchmark tests state and field normalization but has no internal transverse
structure.

\subsection{Benchmark B: two-body scalar spectator}

Let the active constituent have fraction \(x\) and average intrinsic momentum
\(\kT\), with one scalar spectator.  For a normalized scalar wave function
\(\phi(x,\kT)\), the unpolarized canonical overlap is
\begin{equation}
 W^{(2)}(x,\kT,\DT)
 =\cN\,
 \phi^*\!\left(x,\kT+\frac{1-x}{2}\DT\right)
 \phi\!\left(x,\kT-\frac{1-x}{2}\DT\right).
\label{eq:two-body-benchmark}
\end{equation}
This gives immediate tests of:
\begin{align}
 W^{(2)}(x,\kT,\DT)^*&=W^{(2)}(x,\kT,-\DT),
\label{eq:two-body-hermiticity}\\
 W^{(2)}(x,\kT,0)&=\cN|\phi(x,\kT)|^2.
\label{eq:two-body-forward}
\end{align}
For a Gaussian transverse wave function,
\begin{equation}
 \phi(x,\kT)
 =N\,f(x)\exp\!\left[-\frac{\kT^2}{2\beta^2x(1-x)}\right],
\label{eq:gaussian-benchmark}
\end{equation}
the \(\DT\) dependence and Wigner transform are analytic.  This Gaussian is
a verification model, not a physical universal TMD width.

\subsection{Benchmark C: spinor active constituent}

Replace the scalar active kernel by
\crefrange{eq:quark-unpol-kernel}{eq:quark-transverse-kernel} and choose
wave-function components with \(L_z=0,\pm1\).  The benchmark must reproduce:
\begin{itemize}
\item rank-zero unpolarized and helicity projections;
\item rank-one helicity--orbital interference;
\item vanishing of every phase-odd projection for real amplitudes;
\item projector reconstruction of the full \(4\times4\) helicity matrix.
\end{itemize}

\subsection{Benchmark D: three-quark color singlet}

Use a normalized antisymmetric color tensor
\begin{equation}
 C_{abc}=\frac{1}{\sqrt6}\epsilon_{abc}
\label{eq:three-quark-color}
\end{equation}
and a permutation-consistent spin-flavor-momentum amplitude.  Explicit active
summation over the two \(u\) slots and one \(d\) slot must reproduce
\begin{equation}
 N_u=2,
 \qquad
 N_d=1
\label{eq:proton-valence-count}
\end{equation}
for a proton benchmark, with the neutron obtained through a represented
isospin transformation rather than a flavor-equality assignment.

\subsection{Benchmark E: minimal sea and gluon sectors}

Extend the state separately by
\begin{equation}
 |qqq\rangle\oplus|qqqq\bar q\rangle
\label{eq:minimal-sea-benchmark}
\end{equation}
and
\begin{equation}
 |qqq\rangle\oplus|qqqg\rangle.
\label{eq:minimal-gluon-benchmark}
\end{equation}
The first extension must activate positive-\(x\) antiquark overlaps while the
second activates the gluon parent.  Turning the corresponding sector
probability continuously to zero must return the exact zero tests of
\cref{prop:sea-zero-test,prop:gluon-zero-test}.

\subsection{Benchmark F: induced-operator equivalence}

Integrate out the explicit higher sector in a controlled finite-dimensional
model and construct its Feshbach-induced operator.  The full-space matrix
element and reduced-space induced-operator matrix element must agree to the
stated expansion order.  Enabling both simultaneously must trigger the
provenance/double-counting guard.

\section{Contract to Volume III: dynamical Wilson lines}
\label{sec:volume-iii-contract}

Volume~II supplies the zero-rescattering parent
\begin{equation}
 W_{a/N}^{(0)}
 =\langle N|\cO_a\,U_\gamma^{(0)}|N\rangle,
 \qquad
 U_\gamma^{(0)}=\bbone,
\label{eq:zero-rescattering-parent}
\end{equation}
while retaining the full geometric path identity.  Volume~III replaces this
with a controlled expansion or nonperturbative transport,
\begin{equation}
 U_\gamma
 =\bbone-ig\int_\gamma A
 -g^2\int_{\gamma,\mathrm{ord}}AA+\cdots,
\label{eq:wilson-expansion-contract}
\end{equation}
and must return:
\begin{enumerate}
\item explicit intermediate-state denominators and pole prescriptions;
\item a rapidity regulator and soft-overlap subtraction compatible with
      Volume~V;
\item future/past reversal at operator level;
\item flavor-resolved quark phases;
\item ordered gluon links and independent \(f/d\) color contractions;
\item stability with increasing eikonal/Wilson order;
\item a zero-coupling limit equal to \cref{eq:zero-rescattering-parent}.
\end{enumerate}
No scalar phase may be multiplied onto all operators as a substitute.

\section{Contract to Volume IV: spin-1 nuclear matching}
\label{sec:volume-iv-contract}

The nucleon object exported to the nuclear theory is
\begin{equation}
 \mathbb W_N=
 \left\{
 W_{a/N,\lambda'\lambda}^{[J]}
 (x,\kT,\DNT;\mathfrak o,\mathfrak m)
 \right\}_{a,J,\lambda',\lambda},
\label{eq:nucleon-export-object}
\end{equation}
with explicit proton/neutron, flavor, link/color, rank, scheme, and member
identity.  Volume~IV must not consume only named scalar TMD curves when the
nuclear spin recoupling requires the parent helicity matrix.

The nuclear adapter supplies a separate recoil map between the external
hadron transfer \(\DT\) and active-nucleon transfer \(\DNT\).  That relation
must not be inserted into the nucleon overlap itself.  This preserves the
distinction between internal parton recoil and nuclear constituent recoil.

\section{Acceptance criteria for Volume II}
\label{sec:acceptance}

A Volume~II implementation is accepted only when all of the following hold:
\begin{enumerate}
\item The input state satisfies the complete Volume~I state and current
      acceptance gates.
\item Quark, antiquark, and gluon parents are generated from one microscopic
      member and one normalization ledger.
\item The symmetric \(\xi=0\) recoil map is the single implementation
      authority and passes closure, inverse, Jacobian, and physical-momentum
      tests.
\item Every operator acts between explicit incoming and outgoing momentum
      fibers.
\item Operator identity includes path, representation, ordered gluon links,
      color status, regulator, soft scheme, scales, rank convention, and
      truncation.
\item The canonical positive-\(x\) one-body overlap is diagonal in Fock number;
      every off-diagonal block has a named physical source.
\item Quark projectors reconstruct all sixteen twist-two GTMD scalar degrees
      of freedom at non-degenerate kinematics.
\item Antiquarks are computed from explicit positive-\(x\) active constituents
      or a named induced operator, not copied from quarks.
\item Gluons are computed from explicit gluon sectors or a named induced
      operator, with ordered link pairs and no implicit \(f/d\) mixture.
\item Hermiticity, parity adapters, link reversal, and zero-phase limits pass.
\item The forward amplitude correlator is PSD without clipping where the Gram
      theorem applies.
\item TMD, Wigner, GPD, PDF, current, form-factor, OAM, and spin--orbit routes
      are typed reductions of the common parent.
\item Collinear and local-current reductions include explicit matching when
      the staple/soft operator differs from the straight-link operator.
\item Rank-aware Fourier transforms pass round-trip and angular-zero tests.
\item Flavor number, electric charge, helicity, momentum, and angular-momentum
      ledgers close within separately reported truncation and matching errors.
\item Results are stable under declared basis, regulator, Fock-sector, OAM,
      quadrature, and tensor-network refinements.
\item At least one nucleon observable family excluded from calibration is
      predicted from the shared state.
\item Removing every present fitted named-TMD input leaves the microscopic
      overlap calculation defined.
\item All provenance and no-double-counting gates pass.
\item The exported nucleon object is sufficient for Volume~III Wilson dynamics
      and Volume~IV spin-resolved nuclear matching without reconstructing lost
      flavor, helicity, color, or operator information.
\end{enumerate}

Only after these conditions are met may the parent be called a predictive
microscopic nucleon boundary.  Completion of the typed C1 migration alone is
necessary infrastructure, not completion of this dynamical volume.

\appendix

\section{Matrix form of the zero-skewness recoil map}
\label{app:recoil-matrix}

For one active index \(j\), collect the independent intrinsic coordinates
into a vector \(K\).  The incoming/outgoing coordinates are affine maps
\begin{equation}
 K_{\rm in}=K+A_j(x)\DT,
 \qquad
 K_{\rm out}=K-A_j(x)\DT,
\label{eq:recoil-affine-matrix}
\end{equation}
where the active row of \(A_j\) is \(-(1-x_j)/2\) and each spectator row is
\(+x_i/2\), after choosing independent coordinates compatible with the
constraint.  The column sum is zero.  Therefore
\begin{equation}
 \det\frac{\partial K_{\rm in}}{\partial K}
 =\det\frac{\partial K_{\rm out}}{\partial K}=1.
\label{eq:recoil-matrix-jacobian}
\end{equation}
The average and difference are
\begin{equation}
 \frac{K_{\rm out}+K_{\rm in}}2=K,
 \qquad
 K_{\rm out}-K_{\rm in}=-2A_j(x)\DT.
\label{eq:recoil-average-difference}
\end{equation}
This representation is convenient for automatic differentiation and for
verifying alternative frame conventions.

\section{Good-field spin kernels and charge conjugation}
\label{app:good-field}

\subsection{Normalized active spin matrices}

Let \(\chi_\lambda\) be orthonormal two-component light-front helicity
spinors.  The normalized active kernels obey
\begin{align}
 \chi_{\lambda'}^\dagger\bbone\chi_\lambda
 &=\delta_{\lambda'\lambda},
\label{eq:good-unpol}\\
 \chi_{\lambda'}^\dagger\sigma_3\chi_\lambda
 &=(2\lambda)\delta_{\lambda'\lambda},
\label{eq:good-helicity}\\
 \chi_{\lambda'}^\dagger\sigma_T^j\chi_\lambda
 &=\text{transverse-spin transition matrix}.
\label{eq:good-transverse}
\end{align}
Whether \(\lambda=\pm\tfrac12\) or signed labels \(\pm1\) are stored is part
of the spin convention.  The adapter must not mix the two normalizations.

\subsection{Charge-conjugation parities}

For a conventional charge-conjugation matrix,
\begin{equation}
 C\gamma^{\mu T}C^{-1}=-\gamma^\mu,
 \qquad
 C\gamma_5^T C^{-1}=\gamma_5.
\label{eq:c-matrix-identities}
\end{equation}
Thus the local bilinear parities include
\begin{longtable}{L{0.27\textwidth}L{0.18\textwidth}L{0.45\textwidth}}
\toprule
Bilinear & \(\chi_\Gamma\) & Adapter warning\\
\midrule
\endhead
\(\gamma^\mu\) & \(-1\) & Vector antiquark number enters the local charge with the opposite sign.\\
\(\gamma^\mu\gamma_5\) & \(+1\) & The negative-\(x\) helicity relation differs from the vector relation.\\
\(i\sigma^{\mu\nu}\gamma_5\) & \(-1\) in the standard convention & Endpoint order, Fourier sign, and transverse phase still determine the complete GTMD adapter.\\
\bottomrule
\end{longtable}
These local identities are necessary but not sufficient for the bilocal
negative-\(x\) conversion, which also transforms the Wilson path and external
helicity phases.

\section{Nucleon leading-twist TMD registry}
\label{app:tmd-registry}

\subsection{Quark and antiquark registry}

The table gives a standard spin-\(\tfrac12\) inventory.  Exact prefactors are
stored in the projector adapter.

\begingroup
\setlength{\tabcolsep}{3pt}
\begin{longtable}{L{0.10\textwidth}L{0.16\textwidth}L{0.22\textwidth}L{0.08\textwidth}L{0.11\textwidth}L{0.22\textwidth}}
\toprule
Target & Function & Leading operator content & Rank & Naive \(T\) & Collinear status\\
\midrule
\endhead
U & \(f_1^q\) & \(\gamma^+\), unpolarized quark & 0 & even & direct PDF\\
U & \(h_1^{\perp q}\) & \(i\sigma^{j+}\gamma_5\), transverse quark & 1 & odd & moment only\\
L & \(g_{1L}^q\) & \(\gamma^+\gamma_5\), quark helicity & 0 & even & direct helicity PDF\\
L & \(h_{1L}^{\perp q}\) & \(i\sigma^{j+}\gamma_5\), transverse quark & 1 & even & moment only\\
T & \(f_{1T}^{\perp q}\) & \(\gamma^+\), transverse target & 1 & odd & moment only\\
T & \(g_{1T}^q\) & \(\gamma^+\gamma_5\), transverse target & 1 & even & moment only\\
T & \(h_1^q\) & \(i\sigma^{j+}\gamma_5\), aligned transverse spins & 0 & even & direct transversity PDF\\
T & \(h_{1T}^{\perp q}\) & \(i\sigma^{j+}\gamma_5\), rank-two spin tensor & 2 & even & weighted moment\\
\bottomrule
\end{longtable}
\endgroup
The symbol \(h_1\) may denote a particular combination of
\(h_{1T}\) and \(h_{1T}^{\perp}\) depending on convention.  The registry
stores that linear adapter explicitly.  The same rows apply to positive-
\(x\) antiquarks with independent flavor and charge-conjugation metadata.

\subsection{Gluon registry}

\begingroup
\setlength{\tabcolsep}{3pt}
\begin{longtable}{L{0.10\textwidth}L{0.17\textwidth}L{0.22\textwidth}L{0.08\textwidth}L{0.11\textwidth}L{0.21\textwidth}}
\toprule
Target & Function & Gluon polarization tensor & Rank & Naive \(T\) & Link/color requirement\\
\midrule
\endhead
U & \(f_1^g\) & trace & 0 & even & ordered pair stored; color status may be unresolved in zero-phase core\\
U & \(h_1^{\perp g}\) & symmetric traceless & 2 & even & ordered pair stored\\
L & \(g_{1L}^g\) & antisymmetric/helicity & 0 & even & ordered pair stored\\
L & \(h_{1L}^{\perp g}\) & symmetric traceless & 2 & odd & explicit future/past and \(f/d\) status\\
T & \(f_{1T}^{\perp g}\) & trace with transverse target & 1 & odd & explicit future/past and \(f/d\) status\\
T & \(g_{1T}^g\) & antisymmetric with transverse target & 1 & even & ordered pair stored\\
T & \(h_{1T}^{g}\) & symmetric traceless & 1 & odd & explicit future/past and \(f/d\) status\\
T & \(h_{1T}^{\perp g}\) & symmetric traceless & 3 & odd & explicit future/past and \(f/d\) status\\
\bottomrule
\end{longtable}
\endgroup
The rank is the irreducible transverse weight of the complete modulation, not
merely the number of free indices on the field-strength correlator.

\section{Standard quark-GTMD basis adapter}
\label{app:mms-adapter}

A commonly used twist-two convention writes the vector and axial sectors as
\cite{Meissner2009}
\begin{align}
 W^{[\gamma^+]}
 &=\frac{1}{2M_N}\bar u(p',\Lambda')
 \left[
 F_{1,1}
 +\frac{i\sigma^{i+}k_T^i}{P^+}F_{1,2}
 +\frac{i\sigma^{i+}\Delta_T^i}{P^+}F_{1,3}
 +\frac{i\sigma^{ij}k_T^i\Delta_T^j}{M_N^2}F_{1,4}
 \right]u(p,\Lambda),
\label{eq:mms-vector-adapter}\\
 W^{[\gamma^+\gamma_5]}
 &=\frac{1}{2M_N}\bar u(p',\Lambda')
 \left[
 -\frac{i\epsilon_T^{ij}k_T^i\Delta_T^j}{M_N^2}G_{1,1}
 +\frac{i\sigma^{i+}\gamma_5 k_T^i}{P^+}G_{1,2}
 +\frac{i\sigma^{i+}\gamma_5\Delta_T^i}{P^+}G_{1,3}
 +i\sigma^{+-}\gamma_5G_{1,4}
 \right]u(p,\Lambda).
\label{eq:mms-axial-adapter}
\end{align}
A corresponding chiral-odd basis is
\begin{align}
 W^{[i\sigma^{j+}\gamma_5]}
 &=\frac{1}{2M_N}\bar u(p',\Lambda')
 \Bigg[
 -\frac{i\epsilon_T^{ij}k_T^i}{M_N}H_{1,1}
 -\frac{i\epsilon_T^{ij}\Delta_T^i}{M_N}H_{1,2}
 +\frac{M_N i\sigma^{j+}\gamma_5}{P^+}H_{1,3}
\nonumber\\
 &\quad
 +\frac{k_T^j i\sigma^{k+}\gamma_5k_T^k}{M_NP^+}H_{1,4}
 +\frac{\Delta_T^j i\sigma^{k+}\gamma_5k_T^k}{M_NP^+}H_{1,5}
 +\frac{\Delta_T^j i\sigma^{k+}\gamma_5\Delta_T^k}{M_NP^+}H_{1,6}
\nonumber\\
 &\quad
 +\frac{k_T^j i\sigma^{+-}\gamma_5}{M_N}H_{1,7}
 +\frac{\Delta_T^j i\sigma^{+-}\gamma_5}{M_N}H_{1,8}
 \Bigg]u(p,\Lambda).
\label{eq:mms-tensor-adapter}
\end{align}
Equations
\crefrange{eq:mms-vector-adapter}{eq:mms-tensor-adapter} define one
reference adapter, not the internal data model.  The implementation verifies
the equations with explicit spinors and compares their Gram-generated dual
projectors with the reference convention.  A different sign for
\(\epsilon_T^{12}\), a different \(\sigma^{\mu\nu}\) definition, or a
different placement of \(2M_N\) requires a distinct adapter ID.

\section{Projector-generation algorithm}
\label{app:projector-algorithm}

For each operator sector:
\begin{enumerate}
\item generate all candidate covariants with the required parity, twist,
      free indices, mass dimension, and transverse weight;
\item reduce them using on-shell spinor identities, Gordon identities,
      transverse Schouten identities, and the selected \(\xi\) restriction;
\item evaluate the candidate matrices at several generic rational or
      high-precision kinematic points;
\item compute the numerical/symbolic rank and select a stable independent
      basis;
\item form the Gram matrix \cref{eq:gtmd-gram};
\item construct dual projectors and verify
      \(\langle\cP_A,B_B\rangle=\delta_{AB}\);
\item compare with every installed literature adapter through an invertible
      linear transformation;
\item repeat at degenerate limits and construct reduced projector bases for
      TMD, GPD, and local-current spaces.
\end{enumerate}
The generic-point condition should avoid
\begin{equation}
 \kT=0,
 \qquad
 \DT=0,
 \qquad
 \kT\parallel\DT,
\label{eq:projector-degenerate-points}
\end{equation}
because those loci reduce the rank of the complete GTMD basis.

\section{Nonzero-skewness extension contract}
\label{app:nonzero-xi}

For \(\xi\ne0\),
\begin{equation}
 p^+=(1+\xi)P^+,
 \qquad
 p'^+=(1-\xi)P^+.
\label{eq:nonzero-xi-hadron}
\end{equation}
In the DGLAP quark region \(x>\xi\), a standard average-variable recoil map
has active fractions
\begin{equation}
 x_j^{\rm in}=\frac{x_j+\xi}{1+\xi},
 \qquad
 x_j^{\rm out}=\frac{x_j-\xi}{1-\xi},
\label{eq:nonzero-xi-active-fractions}
\end{equation}
and spectator fractions
\begin{equation}
 x_i^{\rm in}=\frac{x_i}{1+\xi},
 \qquad
 x_i^{\rm out}=\frac{x_i}{1-\xi}.
\label{eq:nonzero-xi-spectator-fractions}
\end{equation}
A compatible symmetric transverse mapping is
\begin{align}
 \bm\kappa_{jT}^{\rm in}
 &=\bm{k}_{jT}-\frac{1-x_j}{1+\xi}\frac{\DT}{2},
 &\bm\kappa_{jT}^{\rm out}
 &=\bm{k}_{jT}+\frac{1-x_j}{1-\xi}\frac{\DT}{2},
\label{eq:nonzero-xi-active-recoil}\\
 \bm\kappa_{iT}^{\rm in}
 &=\bm{k}_{iT}+\frac{x_i}{1+\xi}\frac{\DT}{2},
 &\bm\kappa_{iT}^{\rm out}
 &=\bm{k}_{iT}-\frac{x_i}{1-\xi}\frac{\DT}{2}.
\label{eq:nonzero-xi-spectator-recoil}
\end{align}
The exact Jacobian and state-normalization powers depend on the Fock-state
convention and must be derived rather than copied.  At \(\xi=0\), these
expressions reduce to \cref{def:recoil-map}.

In the ERBL region \(|x|<\xi\), the leading overlap contains
\(n+1\leftrightarrow n-1\) Fock sectors.  A complete extension must add:
\begin{enumerate}
\item positive- and negative-\(x\) support with exact charge conjugation;
\item DGLAP and ERBL sector maps with compatible endpoint behavior;
\item local-moment polynomiality through the highest retained moment;
\item continuity constraints at \(x=\pm\xi\);
\item pair-creation operator blocks and their counterterms;
\item convergence under enlargement of the sea and gluon sectors.
\end{enumerate}
The present zero-skewness volume must not be advertised as a complete
polynomiality construction.

\section{Machine-readable operator and closure sketches}
\label{app:machine-readable}

A quark operator identity may be serialized as
\begin{verbatim}
{
  "schema": "deuteronwigner.gtmd_operator.v1",
  "species": "quark",
  "flavor": "u",
  "projection": "gamma_plus",
  "incoming_fiber": "N:Pplus:DeltaT_minus_half:r17",
  "outgoing_fiber": "N:Pplus:DeltaT_plus_half:r17",
  "xi": 0.0,
  "wilson_path": "staple_future_fundamental_v1",
  "uv_regulator": "LF_basis_r17",
  "rapidity_regulator": "UNRESOLVED",
  "soft_subtraction": "UNRESOLVED",
  "mu": null,
  "zeta": null,
  "scheme": "regulated_LF_overlap",
  "rank_spec": "vector_GTMD_basis_MMS_v1",
  "microscopic_member": "member_0042"
}
\end{verbatim}
Such an object is valid for a regulated overlap.  It must fail closed if
passed directly to a process requiring a subtracted TMD scheme.

A closure report may contain
\begin{verbatim}
{
  "parent_id": "member_0042:u:gamma_plus:xi0",
  "recoil_id": "SYMMETRIC_XI0_v1",
  "tests": {
    "hermiticity_max_abs": 2.1e-13,
    "forward_limit_max_abs": 8.0e-14,
    "wigner_tmd_marginal_rel": 4.2e-10,
    "gpd_current_route_rel": 7.3e-4,
    "forward_psd_min_eigenvalue": 1.7e-8
  },
  "error_budget": {
    "basis_truncation": 5.0e-4,
    "fock_truncation": 1.2e-3,
    "matching": 6.0e-4,
    "quadrature": 2.0e-7
  },
  "accepted": true
}
\end{verbatim}
The numbers are illustrative schema values, not project acceptance results.

\section{Formalism-to-computation work-package map}
\label{app:work-packages}

The recommended implementation sequence after the typed C1 spine is:
\begin{longtable}{L{0.13\textwidth}L{0.31\textwidth}L{0.46\textwidth}}
\toprule
Package & Scientific object & Completion test\\
\midrule
\endhead
C2A & \texttt{MomentumFiber} and \texttt{RecoilMap} & Exact analytic benchmarks A--D and mismatch injection.\\
C2B & Common diagonal \texttt{OverlapKernel} & Quark vector-current number and Hermiticity closure on toy Fock states.\\
C2C & Quark helicity matrix and projector engine & Sixteen-function rank at generic kinematics and reduced-rank TMD/GPD limits.\\
C2D & Positive-\(x\) antiquark and gluon active selectors & Exact sea/gluon zero tests and activation with minimal higher sectors.\\
C2E & Typed reduction maps & TMD/Wigner/GPD/PDF/current route closure on analytic benchmarks.\\
C2F & Common-parent closure and convergence manifests & Machine-readable route residuals and independent truncation axes.\\
C3 & Dynamical Wilson-line engine & Volume~III acceptance, including phases and \(f/d\) color.\\
\bottomrule
\end{longtable}
The exact Codex prompt should be written after the C1 implementation report so
that it uses the actual repository APIs, file paths, and unresolved adapter
fields.

\small
\begin{thebibliography}{99}

\bibitem{Dirac1949}
P.~A.~M.~Dirac,
``Forms of relativistic dynamics,''
Rev. Mod. Phys. \textbf{21}, 392 (1949).

\bibitem{BrodskyPauliPinsky1998}
S.~J.~Brodsky, H.-C.~Pauli, and S.~S.~Pinsky,
``Quantum chromodynamics and other field theories on the light cone,''
Phys. Rept. \textbf{301}, 299 (1998),
arXiv:hep-ph/9705477.

\bibitem{Diehl2000}
M.~Diehl, T.~Feldmann, R.~Jakob, and P.~Kroll,
``The overlap representation of skewed quark and gluon distributions,''
Nucl. Phys. B \textbf{596}, 33 (2001), Erratum B \textbf{605}, 647 (2001),
arXiv:hep-ph/0009255.

\bibitem{Diehl2003}
M.~Diehl,
``Generalized parton distributions,''
Phys. Rept. \textbf{388}, 41 (2003),
arXiv:hep-ph/0307382.

\bibitem{Meissner2009}
S.~Mei{\ss}ner, A.~Metz, and M.~Schlegel,
``Generalized parton correlation functions for a spin-1/2 hadron,''
JHEP \textbf{08}, 056 (2009),
arXiv:0906.5323.

\bibitem{BelitskyJiYuan2004}
A.~V.~Belitsky, X.~Ji, and F.~Yuan,
``Quark imaging in the proton via quantum phase-space distributions,''
Phys. Rev. D \textbf{69}, 074014 (2004),
arXiv:hep-ph/0307383.

\bibitem{LorcePasquini2011}
C.~Lorc\'e and B.~Pasquini,
``Quark Wigner distributions and orbital angular momentum,''
Phys. Rev. D \textbf{84}, 014015 (2011),
arXiv:1106.0139.

\bibitem{Lorce2012}
C.~Lorc\'e,
``Wilson lines and orbital angular momentum,''
Phys. Lett. B \textbf{719}, 185 (2013),
arXiv:1210.2581.

\bibitem{LorcePasquini2013}
C.~Lorc\'e and B.~Pasquini,
``Structure analysis of the generalized correlator of quark and gluon for a
spin-1/2 target,''
JHEP \textbf{09}, 138 (2013),
arXiv:1307.4497.

\bibitem{Kanazawa2014}
K.~Kanazawa, C.~Lorc\'e, A.~Metz, B.~Pasquini, and M.~Schlegel,
``Twist-2 generalized transverse-momentum dependent parton distributions and
the spin/orbital structure of the nucleon,''
Phys. Rev. D \textbf{90}, 014028 (2014),
arXiv:1403.5226.

\bibitem{Ji1997}
X.~Ji,
``Gauge-invariant decomposition of nucleon spin,''
Phys. Rev. Lett. \textbf{78}, 610 (1997),
arXiv:hep-ph/9603249.

\bibitem{Hatta2011}
Y.~Hatta,
``Gluon polarization in the nucleon demystified,''
Phys. Rev. D \textbf{84}, 041701 (2011),
arXiv:1101.5989.

\bibitem{MuldersRodrigues2001}
P.~J.~Mulders and J.~Rodrigues,
``Transverse momentum dependence in gluon distribution and fragmentation
functions,''
Phys. Rev. D \textbf{63}, 094021 (2001),
arXiv:hep-ph/0009343.

\bibitem{Buffing2013}
M.~G.~A.~Buffing, A.~Mukherjee, and P.~J.~Mulders,
``Generalized universality of definite rank gluon transverse momentum
dependent correlators,''
Phys. Rev. D \textbf{88}, 054027 (2013),
arXiv:1306.5897.

\bibitem{Boer2016}
D.~Boer, S.~Cotogno, T.~van Daal, P.~J.~Mulders, A.~Signori, and
Y.-J.~Zhou,
``Gluon and Wilson loop TMDs for hadrons of spin \(\le1\),''
JHEP \textbf{10}, 013 (2016),
arXiv:1607.01654.

\bibitem{Cotogno2017}
S.~Cotogno, T.~van Daal, and P.~J.~Mulders,
``Positivity bounds on gluon TMDs for hadrons of spin \(\le1\),''
JHEP \textbf{11}, 185 (2017),
arXiv:1709.07827.

\bibitem{Collins2011}
J.~Collins,
\emph{Foundations of Perturbative QCD},
Cambridge University Press (2011).

\bibitem{Echevarria2016}
M.~G.~Echevarria, A.~Idilbi, K.~Kanazawa, C.~Lorc\'e, A.~Metz,
B.~Pasquini, and M.~Schlegel,
``Proper definition and evolution of generalized transverse momentum
dependent distributions,''
Phys. Lett. B \textbf{759}, 336 (2016),
arXiv:1602.06953.

\bibitem{Bacchetta2007}
A.~Bacchetta, M.~Diehl, K.~Goeke, A.~Metz, P.~J.~Mulders, and M.~Schlegel,
``Semi-inclusive deep inelastic scattering at small transverse momentum,''
JHEP \textbf{02}, 093 (2007),
arXiv:hep-ph/0611265.

\bibitem{BomhofMuldersPijlman2006}
C.~J.~Bomhof, P.~J.~Mulders, and F.~Pijlman,
``Gauge link structure in quark-quark correlators in hard processes,''
Phys. Lett. B \textbf{596}, 277 (2004),
arXiv:hep-ph/0406099;
and related process-link analyses.

\end{thebibliography}

\end{document}
